\chapter{Action profiles and universal recipients}
\label{chap:vol2-recipients}

A universal boundary action must receive actions together with their
comparison homotopies. A multiplication on the underlying coefficient
algebra does not suffice. For an action whose first nonzero term is
quadratic in its parameter, doubling that parameter multiplies the term
by four, whereas composing two copies adds their first terms and
multiplies it by two. The standard additive multiplication is therefore
incompatible with such an action whenever its first class is nonzero.
Nevertheless, a nonlinear change of actor can place the action inside
a multiplicative one.

The higher terms of these actions have a simple scalar example. Over a
characteristic-zero field, the formal solution of
\[
 p=q+\frac{z^2}{2}p^2,\qquad p\equiv q\pmod z,
\]
is
\[
 p=\sum_{n\geq0}\frac{1}{n+1}\binom{2n}{n}
                 \left(\frac{z^2}{2}\right)^n q^{n+1}.
\]
Writing $p=qC(z^2q/2)$ reduces the equation to $C=1+uC^2$,
whose solution with constant term one is
$(1-\sqrt{1-4u})/(2u)$. For noncommuting coefficients the same
implicit construction uses ordered trees. Its terms define actual
actions on the collision coefficients. The universal question then
concerns all maps between these actions, including their higher paths.

\begingroup
\newcommand{\vCE}{\operatorname{CE}_{\star}}
\newcommand{\vU}{U_{\star}}
\newcommand{\vB}{\operatorname{B}_{\star}}
\newcommand{\vSym}{\operatorname{Sym}_{\star}}
\newcommand{\vRH}{\operatorname{RHom}}
\newcommand{\vMap}{\operatorname{Map}}
\newcommand{\vot}{\mathbin{\otimes^!}}
\newcommand{\vM}{\mathscr M}
\newcommand{\vF}{\mathscr F}
\newcommand{\vAct}{\mathscr A_C}
\newcommand{\vMon}{\mathscr B_C}
\newcommand{\vEnd}{\mathcal E_C}
\newcommand{\vid}{\mathrm{id}}

\section{The coefficients of an action}
\label{sec:vol2-action-carrier}

Fix a field $k$ of characteristic zero and $X=\mathbb A^1_k$.
At a collision, the order of the coinciding labels is retained.
For a nonempty finite set $I$, take a copy $X^I_\sigma$ of $X^I$
for each total order $\sigma$ on $I$. Glue two copies by the identity
on the open
\[
 V_{\sigma\rho}=X^I\setminus
 \bigcup_{\substack{\{i,j\}\subset I\\
                  \sigma,\rho\text{ order }i,j\text{ oppositely}}}
                  \{z_i=z_j\}.
\]
The resulting scheme, denoted by $\mathscr R_I$, is the ordered
configuration scheme of \cite{AKY2024}*{\S4.1}. The gluing condition
holds because, whenever two orders disagree on a pair, at least one
disagrees with any third order on that pair. Thus
$V_{\sigma\rho}\cap V_{\sigma\eta}
 =V_{\sigma\rho}\cap V_{\rho\eta}$.
The scheme is smooth, with a finite affine-chart cover and
quasi-compact intersections. It can be nonseparated. Over a tuple with
equality blocks $B$, the map $\mathscr R_I\to X^I$ has
$\prod_B|B|!$ points, one for each order inside each block.

A surjection $\alpha:I\twoheadrightarrow J$ equipped with an order
on each fibre defines
\[
 \delta_\alpha:\mathscr R_J\longrightarrow\mathscr R_I,
 \qquad z_i=z_{\alpha(i)}.
\]
On charts it expands each label into its ordered fibre. Fibre orders
compose by concatenation. The chart maps glue, since the inverse image
of an overlap is the overlap of the corresponding orders on $J$.
Their compositions are the maps for the composite ordered surjections.
In particular, a nonconsecutive collision uses its actual map into
$\mathscr R_I$; it is not replaced by the small diagonal of one chart.

Let $\mathscr C$ be the enhanced category of compatible complexes of
right differential modules on these schemes:
\begin{equation}\label{eq:vol2-action-category}
 \mathscr C=D(k)\times
     \varprojlim_{I,\alpha}D\!\operatorname{Mod}(\mathscr R_I),
 \qquad \delta_\alpha^!F_I\simeq F_J.
\end{equation}
Differential modules on $\mathscr R_I$ are defined by descent on its
smooth chart cover. The displayed equivalences include all composition
homotopies. The first factor is the isolated empty coefficient
$F_\varnothing$. An object is \emph{positive} if this coefficient is
zero. Complexes are cohomological, suspension $s=[1]$ has degree $-1$,
and every permutation uses the Koszul sign. Extraordinary inverse
images and the diagonal tensor retain their canonical density lines
and shifts. Passing to left modules for a coordinate computation uses
the canonical-bundle equivalence.

For an ordered decomposition $I=A\sqcup B$, let $U_{A,B}$ be the open
where every point in $A$ is distinct from every point in $B$.
Collisions inside either part remain. Restriction of orders identifies
this open with the disjointness open in
$\mathscr R_A\times\mathscr R_B$. With $j_{A,B}$ its inclusion, set
\begin{equation}\label{eq:vol2-action-convolution}
 (F\star G)_I=
 \bigoplus_{I=A\sqcup B}
 j_{A,B,*}\bigl((F_A\boxtimes G_B)|_{U_{A,B}}\bigr).
\end{equation}
An empty part contributes scalar tensor product. The unit $K$ has
$K_\varnothing=k$ and zero nonempty coefficients. Iterating
\eqref{eq:vol2-action-convolution} sums over ordered partitions into
the required number of parts. All parenthesizations impose the same
cross-part inequalities. This proves associativity and its coherences;
permuting the parts gives symmetry. Under a collision, a summand
vanishes if one collision fibre meets two parts. Every surviving
partition descends uniquely to the quotient labels. Localization and
closed-support base change therefore give the collision coherences.
These operations preserve colimits separately.

The second tensor product uses coefficients at the same configuration:
\[
 (F\vot G)_I=\Delta_{\mathscr R_I}^!(F_I\boxtimes G_I).
\]
Its unit $J$ is the family of dualizing objects
$J_I=(\mathscr R_I\to\operatorname{Spec}k)^!k$, with empty
coefficient $k$. The objects $J$ and $K$ have different nonempty
coefficients. The forward interchange
\begin{equation}\label{eq:vol2-action-interchange}
 (F\star G)\vot(H\star L)\longrightarrow
       (F\vot H)\star(G\vot L)
\end{equation}
keeps summands with the same ordered partition and applies the
localization comparison there. In a rectangular iterated product,
either order of interchange keeps exactly the common partitions.
This proves its coherence and gives a diagonal tensor on
coassociative $\star$-coalgebras.

Write $\vM$ for connected coaugmented counital coherent
coassociative coalgebras in $\mathscr C$. Connected means that the
empty coefficient and its unit maps are $k$. A coalgebra $F$ is
\emph{factorizing} if every ordered partition
$P=(I_1|\cdots|I_r)$ into nonempty parts has an equivalence
\begin{equation}\label{eq:vol2-action-factorization}
 \delta_{P,F}:j_P^!F_I\xrightarrow{\sim}
       (F_{I_1}\boxtimes\cdots\boxtimes F_{I_r})|_{U_P}.
\end{equation}
The coproduct supplies this map and its refinement homotopies.
Let $\vF\subset\vM$ be the full subcategory of such coalgebras.
The diagonal tensor preserves it by
\eqref{eq:vol2-action-interchange}.

Fix a boundary $C\in\vF$. A \emph{pointed action} consists of
$F\in\vF$, a point $u:J\to F$, an action
$\alpha:F\vot C\to C$ in $\vM$, and a specified path
$\alpha(u\vot\vid_C)\simeq\vid_C$.
A morphism $(F,\alpha)\to(G,\beta)$ includes a pointed coalgebra
map $f:F\to G$, a path
$\alpha\simeq\beta(f\vot\vid_C)$, and its compatibility on $J$.
Their higher homotopies are retained. Denote this category by $\vAct$.
The functor it classifies is
\begin{equation}\label{eq:vol2-action-functor}
 \Phi_C(F,u)=\vMap^{\vM}_{C/}(F\vot C,C),
\end{equation}
the homotopy fibre of restriction over the identity of $C$.
Representability of $\Phi_C$ is equivalent to a final object of
$\vAct$.

The ordered product of two actions is their composition, with actor
$F\vot G$. Denote it by $a\circledast b$; the first action is
outermost. The unit action on $J$ is initial. A \emph{multiplicative
action} has an associative unital monoid actor for $\vot$ and a
coherent module structure on $C$. Its binary path compares acting
after multiplying the two actors with their composed action. All
associativity and unit homotopies are part of the structure.
Write $\vMon$ for the category of multiplicative actions, with monoid
maps and coherent action paths. A local centre, in this categorical
sense, is a final object of $\vMon$. A final pointed action has a
unique coherent multiplicative structure. The converse requires a
comparison with pointed actions whose actors have no specified
multiplication.

\section{Trees integrating a derivation}
\label{sec:vol2-action-profiles}

Let $\mathfrak g$ be a positive Lie object in $\mathscr C$ whose
coefficient modules lie in degree zero, and put $A=\vU\mathfrak g$.
Fix a degree-zero derivation $D$ of the full coefficient family.
Its extension to $A$ satisfies Leibniz and commutes with all
collision and normal-transfer maps. Products below are ordered products in $A$.
PBW symmetrization uses averaging in characteristic zero.

Give an odd symbol $e$ degree one and zero differential, and form
$\mathfrak h=\mathfrak g\ltimes e\mathfrak g$, with
\[
 [x,em]=e[x,m],\qquad [em,en]=0.
\]
The Chevalley coalgebra is
$\vCE(\mathfrak h)=K\oplus\bigoplus_{r\geq1}
\operatorname{Sym}^r_\star(s\mathfrak h)$, with unshuffle coproduct.
Its coderivation has
$Q_1(sx)=-s(dx)$ and
$Q_2(sx\,sy)=(-1)^{|x|}s[x,y]$.
The \emph{actor weight} counts the factors $s(em)$, which have degree
zero. It is independent of PBW output number.

A degree-one twisting cochain to $A$ can be nonzero only on one
ordinary suspended input. Write its value as
$\theta_w(x;m_1,\ldots,m_w)$ and put $\theta_0(x)=x$.
For $m_S=\mathop{\circ}_{i\in S}m_i$ in
$P_{|S|}=\vSym^{|S|}\mathfrak g$, the twisting equation is
\begin{equation}\label{eq:vol2-profile-MC}
\begin{aligned}
0={}&[x,\theta_w(y;m)]-\theta_w(y;[x,m])
       -[y,\theta_w(x;m)]+\theta_w(x;[y,m])
       -\theta_w([x,y];m)\\
 &+\sum_{\substack{S\sqcup T=[w]\\S,T\ne\varnothing}}
 \bigl(\theta_{|S|}(x;m_S)\theta_{|T|}(y;m_T)
       -\theta_{|S|}(y;m_S)\theta_{|T|}(x;m_T)\bigr).
\end{aligned}
\end{equation}
Here $[x,m]$ is the sum of brackets with each module input, and the
splitting $(S,T)$ is ordered. To verify the signs, use convolution
\[
 (f*g)(c)=\sum_{(c)}(-1)^{|g||c'|}f(c')g(c'').
\]
On two ordinary suspended inputs, the tensor sign of two degree-one
cochains is minus. The terms with no actor input on one side give
the four adjoint terms above; the Chevalley bracket gives the fifth.
All other components vanish by degree. Thus
\eqref{eq:vol2-profile-MC} is both necessary and sufficient.

For a nonempty set $B$ of actor labels, define
\[
 u_B=\frac1{|B|!}\sum_{\sigma\in S_B}
                   m_{\sigma(1)}\cdots m_{\sigma(|B|)}.
\]
Let $\mathcal T_r$ consist of nonnegative tuples
$(a,b_1,\ldots,b_r)$ such that
\[
 a+\sum_i b_i=r,\qquad
 a+\sum_{i=1}^j b_i\geq j+1\quad(0\leq j<r).
\]
These are the child counts of plane rooted trees in preorder, with
the root first. For envelope inputs $v_1,\ldots,v_r$, set
\begin{equation}\label{eq:vol2-profile-trees}
 T_r^D(x;v_1,\ldots,v_r)=
 \sum_{(a,b)\in\mathcal T_r}\frac{D^ax}{a!}
 \sum_{\sigma\in S_r}
 \frac{D^{b_1}v_{\sigma(1)}\cdots D^{b_r}v_{\sigma(r)}}
      {b_1!\cdots b_r!},\qquad T_0^D(x)=x.
\end{equation}
The slots are symmetric; each displayed product remains ordered.
For a scalar sequence $\mathbf t=(t_1,t_2,\ldots)$, define
\begin{equation}\label{eq:vol2-profile-partitions}
 \theta_w^{\mathbf t}(x;m_1,\ldots,m_w)=
 \sum_{\pi\in\operatorname{Part}([w])}
       \left(\prod_{B\in\pi}t_{|B|}\right)
       T_{|\pi|}^D(x;u_B:B\in\pi).
\end{equation}
The partitions are unordered, with nonempty blocks. For example,
\begin{align*}
 \theta_1^{\mathbf t}(x;m)&=t_1Dx\,m,\\
 \theta_2^{\mathbf t}(x;m,n)
 &=t_2Dx\,\frac{mn+nm}{2}
   +t_1^2\left(D^2x\,\frac{mn+nm}{2}
                   +Dx(Dm\,n+Dn\,m)\right).
\end{align*}
In weight three the term proportional to $t_1t_2$ is the sum of
$T_2^D(x;m_i,u_B)$ over $\{i\}\sqcup B=[3]$.

\begin{theorem}[Simultaneous profiles]\label{thm:vol2-profiles}
The maps \eqref{eq:vol2-profile-partitions} satisfy
\eqref{eq:vol2-profile-MC}. They define a map of coassociative
coalgebras
\[
 \vCE(\mathfrak g\ltimes e\mathfrak g)
       \longrightarrow\vCE(\mathfrak g)
\]
under $\vCE(\mathfrak g)$, in the derived category of coalgebras.
Every nonconstant term on a configuration $I$ has actor weight
$w<|I|$. Thus each scalar sequence defines coefficient maps without
completing states or normal derivatives.
\end{theorem}

\begin{proof}
Use a central formal variable $z$ and define right substitution by
\[
 f\{R\}=\sum_{j\geq0}\frac{D^jf}{j!}R^j,
 \qquad U=\sum_{r\geq1}\frac{t_rz^r}{r!}m^r,
 \qquad R=U\{R\},\qquad G_x=x\{R\}.
\]
Every coefficient of $R\in zA[[z]]$ is determined by the preceding
ones. Let $\partial_n$ vary $m$ in direction $n$, including its
$D$-derivatives. Adjoint equivariance gives
$\partial_{[x,m]}U=[x,U]$. Put $U_j=D^jU/j!$ and $X_i=D^ix/i!$.
Differentiating the fixed-point equation gives a linear filtered
equation for $\partial_{[x,m]}R$. The same equation is satisfied by
$[G_x,R]$: the difference of its expanded sides is
\[
 \sum_{i,j}X_i[R^i,U_j]R^j
 =\sum_iX_i\sum_{p+q=i-1}R^p
                 \left(\sum_j[R,U_j]R^j\right)R^q=0,
\]
because the inner sum is $[R,R]$. The linear term of the equation
raises the $z$-filtration, so uniqueness gives
$\partial_{[x,m]}R=[G_x,R]$.

Leibniz for $D$ now implies
\begin{equation}\label{eq:vol2-profile-transport}
 [G_x,G_y]-G_{[x,y]}
       -\partial_{[x,m]}G_y+\partial_{[y,m]}G_x=0.
\end{equation}
Indeed, with $Y_j=D^jy/j!$, the first two terms differ by
$\sum_{i,j}X_i[R^i,Y_j]R^j-\sum_{i,j}Y_j[R^j,X_i]R^i$.
The two derivative terms cancel precisely these sums. Expanding the
root and its ordered subtrees gives
\eqref{eq:vol2-profile-trees}; allocating the original inputs to the
subtree slots gives \eqref{eq:vol2-profile-partitions}. Its
normalization is
$G_x=x+\sum_{w\geq1}z^w\theta_w^{\mathbf t}(x;m,\ldots,m)/w!$.
Coefficient extraction in \eqref{eq:vol2-profile-transport}, followed
by polarization, proves the twisting equation.

A twisting cochain defines a map into $\vB A$. The Chevalley--bar
comparison identifies $\vB A$ with $\vCE(\mathfrak g)$.
Explicitly, this comparison sends a symmetric word in suspended Lie
inputs to the sum of its signed bar words, without a factorial divisor.
It preserves the unshuffle and deconcatenation coproducts. PBW reduces
its associated graded to the Koszul comparison for a symmetric algebra.
The usual Koszul homotopy, divided by positive polynomial weight,
contracts its augmentation complex. At $I$, both filtrations stop
at $|I|$, so the comparison is an equivalence. Its inverse in the
pointed derived category gives the asserted map.

Each nonconstant term consumes one positive ordinary input and $w$
positive actor inputs. Hence $w+1\leq|I|$. All operations in each
surviving term are full coefficient maps, so the finite formulas
retain collision and normal-transfer compatibility.
\end{proof}

For $D=\partial_q$, $x=m=q$ and $t_2=1$ with the other parameters
zero, right substitution becomes $f(q+R)$. Thus $p=q+R$ satisfies
the scalar equation at the chapter's opening. The noncommutative
proof explains why the order of the powers of $R$ is fixed.

The same construction supplies mixed homotopies. In the relative
convolution algebra with differential $d_0$ at the scalar action,
write the positive cochain as
\[
 \theta_{\mathbf t}=\sum_rt_ru_r+
             \sum_{p\leq q}t_pt_q\eta_{pq}+\cdots.
\]
Its equation gives $d_0\eta_{pq}=-[u_p,u_q]$ for $p<q$ and
$d_0\eta_{pp}=-u_p^2$. These are compatible primitives because they
are coefficients of one Maurer--Cartan element. More explicitly,
give a free graded Lie generator $e_\alpha$ degree one for every
nonzero finite multi-index $\alpha$, with weight $\sum r\alpha_r$,
and set
\begin{equation}\label{eq:vol2-profile-cells}
 de_\alpha=-\frac12
   \sum_{\substack{\beta+\gamma=\alpha\\\beta,\gamma\ne0}}
                  [e_\beta,e_\gamma].
\end{equation}
This is a differential. In the tensor envelope it is
$de_\alpha=-\sum_{\beta+\gamma=\alpha}e_\beta e_\gamma$;
applying $d$ twice gives opposite sums over ordered three-part
splittings. Coefficient extraction from the profile equation shows
that $e_\alpha\mapsto[\mathbf t^\alpha]\theta_{\mathbf t}$ is a
dg Lie map. Only the positive-weight completion is used.

\section{A change of actor and an actual bar composite}
\label{sec:vol2-actor-composite}

For a vertex algebra $V$, let $L_V$ be its main-diagonal Lie family
on \eqref{eq:vol2-action-category}, with singleton the right
differential module of states and connection induced by translation.
In each ordered chart its higher coefficient is the full normal
direct image of this module on the small diagonal. The chiral bracket
uses principal parts. Put $C_V=\vCE(L_V)$.
For the vacuum vertex algebra $k$, this is the dualizing coalgebra
$J$. For $E=\Lambda(e)$, $|e|=1$, with zero differential and
translation, write $C_E$ for its factorizing coalgebra.

The following tensor comparison specifies how the algebraic profiles
act on the coefficient family.

\begin{lemma}\label{lem:vol2-strong-boundary-tensor}
If $C$ is a cocommutative factorizing coalgebra with empty
coefficient $k$, then $F\mapsto F\vot C$ is a colimit-preserving
strong symmetric monoidal functor for $\star$. For $C=C_V$ it
sends $L_k$ to $L_V$ and $L_E$ to $L_V\ltimes eL_V$, with their
Lie structures. Consequently
\begin{equation}\label{eq:vol2-exterior-comparison}
 C_E\vot C_V\simeq\vCE(L_V\ltimes eL_V).
\end{equation}
The comparison preserves the vacuum, actor grading, finite scalar
extensions, and all collision maps and normal derivatives.
\end{lemma}

\begin{proof}
For each ordered partition, tensor its convolution summand with
$C_I$ and use \eqref{eq:vol2-action-factorization} for $C$.
The projection formula for localization identifies the result with
the corresponding summand of
$(F\vot C)\star(G\vot C)$. At the empty coefficient the map is
the scalar identity. All iterated comparisons use the same ordered
partition and the same iterated coproduct of $C$. This gives
associativity and units; cocommutativity gives symmetry. Collision
pullback keeps precisely the saturated partitions on both sides,
with the same base-change maps. Thus the functor is strong monoidal
on the enhanced diagram and commutes with its Chevalley complexes,
envelopes and bars.

On the main diagonal, the projection formula identifies
$L_k\vot C_V$ with $L_V$. The principal-part boundary of the vacuum
line, tensored with the factorization identification of $C_V$ on the
punctured two-point locus, is the chiral bracket of $V$. This checks
the Lie identity on both ordered branches with the full normal
direct image. The even vacuum and odd $e$ summands give the stated
semidirect product. These coefficient formulas preserve actor
scaling and tensor powers. Finite free scalar extensions preserve
their equivalences, which proves the remaining assertions.
\end{proof}

For $\mathfrak g=L_V$ and this comparison, write
$\alpha_{\mathbf t}:C_E\vot C_V\to C_V$ for the profile action.
The sequence $t_1=1$, $t_r=0$ for $r>1$ gives $\alpha_D$.
The actor maps receiving the other profiles are nonlinear even
though a map between the exterior coefficient algebras is linear
on its degree-one generator.

Let $\mathfrak l=L_k$. In $\vU(\mathfrak l\ltimes e\mathfrak l)$
define
\begin{equation}\label{eq:vol2-marked-PBW}
 \rho_r(m_1,\ldots,m_r)=\frac1{r!}\sum_{\sigma\in S_r}
       m_{\sigma(1)}\cdots m_{\sigma(r-1)}(em_{\sigma(r)}).
\end{equation}
This is the averaged PBW product with one marked odd factor.
In particular its diagonal value is $m^{r-1}em$. The cochain
\[
 \tau_{\mathbf t}(sx)=x,\qquad
 \tau_{\mathbf t}(s(em_1)\cdots s(em_r))=t_r\rho_r(m_1,\ldots,m_r)
\]
has all other reduced components zero.

\begin{proposition}\label{prop:vol2-profile-actor}
The cochain $\tau_{\mathbf t}$ is twisting and defines a map
$\Sigma_{\mathbf t}:C_E\to C_E$ under and over $J$.
Its formal coordinate series is
$\sum_{r\geq1}t_rp^r/r!$. There is a pointed action path
\begin{equation}\label{eq:vol2-profile-comparison}
 \alpha_D(\Sigma_{\mathbf t}\vot\vid_{C_V})
                    \simeq\alpha_{\mathbf t}.
\end{equation}
It preserves every ordered cut, collision, normal transfer and unit
homotopy. In the associative bar its composite corestriction is
exactly \eqref{eq:vol2-profile-partitions}.
\end{proposition}

\begin{proof}
The maps $\rho_r$ are adjoint-equivariant, giving the twisting
equation with one ordinary input. With two ordinary inputs it is
the equation for the canonical cochain. On pure actor inputs,
polarization reduces the equation to the square of
\[
 \sum_{r\geq1}\frac{t_r}{r!}m^{r-1}em.
\]
It is zero because $[m,em]=0$ and $(em)^2=0$. The induced map
to the associative bar is the finite sum of tensor powers of
$s\tau_{\mathbf t}$ applied to the iterated reduced coproduct.
The inverse Chevalley--bar comparison gives $\Sigma_{\mathbf t}$.
Its restriction to $J$ is the identity and its augmentation is
fixed. Removing the marked odd factor in
\eqref{eq:vol2-marked-PBW} gives the scalar coordinate coefficient
$t_r/r!$, with exponential polarization.

To calculate the composite, put
$B=\vU(\mathfrak g\ltimes e\mathfrak g)$ and let $\pi:B\to A$
kill the odd ideal. There is a degree-$-1$ projection derivation
\begin{equation}\label{eq:vol2-actor-contraction}
 q(x)=0,\quad q(em)=m,\quad q(1)=0,\qquad
 q(uv)=q(u)\pi(v)+(-1)^{|u|}\pi(u)q(v).
\end{equation}
It respects the envelope relations: the mixed relation reduces to
the same bracket identity, while a relation with two odd inputs
maps to zero. In $\operatorname{Hom}(\vB B,A)$ let $a$ be the
canonical length-one twisting cochain induced by $\pi$ and $p$
the length-one cochain induced by $q$. They have degrees one and
zero and satisfy
\[
 da+a^2=0,\qquad dp+[a,p]=0.
\]
Apply $D$ on the output and solve
\[
 R=\sum_{j\geq0}\frac{D^jp}{j!}R^j,
 \qquad A_R=\sum_{j\geq0}\frac{D^ja}{j!}R^j
\]
in positive bar length. These equations give a twisting cochain.
For completeness, put $a_j=D^ja/j!$, $p_j=D^jp/j!$.
The differentiated identities are
$da_j=-\sum_{u+v=j}a_ua_v$ and
$dp_j=-\sum_{u+v=j}[a_u,p_v]$. Moreover,
\[
 \sum_{r,s}[a_r,p_s]R^{r+s}
 =[A_R,R]-\sum_s p_s[A_R,R^s],
\]
since their difference is
$\sum_r a_r[R^r,\sum_s p_sR^s]=0$.
The differentiated fixed-point equation therefore has the unique
solution $dR=-[A_R,R]$. Using that each $a_j$ is odd gives
\[
 dA_R=-\sum_{r,s}a_ra_sR^{r+s}
                   +\sum_r a_r[A_R,R^r],\qquad
 A_R^2-\sum_{r,s}a_ra_sR^{r+s}
                   =\sum_r a_r[R^r,A_R].
\]
Thus $dA_R+A_R^2=0$.

The associated bar map $\mathcal F_D:\vB B\to\vB A$ has
length-$(w+1)$ corestriction on $[v_0|\cdots|v_w]$
\begin{equation}\label{eq:vol2-profile-bar-tree}
 \sum_{(a,b)\in\mathcal T_w}
 \frac{D^a\pi(v_0)}{a!}
 \frac{D^{b_1}q(v_1)\cdots D^{b_w}q(v_w)}{b_1!\cdots b_w!}.
\end{equation}
It vanishes unless the first input has exterior degree zero and the
others have exterior degree one. Restriction to the Chevalley coalgebra
is the tree action $\alpha_D$.

Lemma~\ref{lem:vol2-strong-boundary-tensor} carries the actor map
from $\mathfrak l$ to $\mathfrak g$. Its bar representative
partitions the actor inputs into marked blocks and places the one
ordinary input in a singleton. Formula
\eqref{eq:vol2-profile-bar-tree} forces this singleton first.
Equation~\eqref{eq:vol2-actor-contraction} gives $q\rho_r=u_{[r]}$.
For each unordered partition, the sum over orders of its blocks
is therefore $T_{|\pi|}^D$, with coefficient
$\prod_{B\in\pi}t_{|B|}$. This proves equality of the actual
corestrictions. The pointed Chevalley--bar inverses and the strong
tensor comparison transport that equality to
\eqref{eq:vol2-profile-comparison}, including all unit paths.
Each calculation is finite on every configuration.
\end{proof}

The quadratic map with $t_2=1$ has coordinate series $p^2/2$.
The map with coordinate series $p^2$ has $t_2=2$. This normalization
is independent of the unnormalized permutation sum in the
Chevalley--bar comparison.

\section{Natural homotopies and multiplication}
\label{sec:vol2-natural-actions}

The profile comparison gives a map to $\alpha_D$. To receive
compositions, that target must carry its own coherent multiplication.
The obstruction groups can be computed before evaluating on a
particular boundary.

Let $W$ be finite-dimensional in degree zero and set
$E_W=\Lambda(W[-1])$, whose generators have degree one.
Let $\mathfrak d$ be a Lie algebra of degree-zero derivations of an
ordinary Lie algebra $\mathfrak g$. A \emph{natural operation} uses
each Lie input once, applies elements of $U(\mathfrak d)$ to these
inputs, and multiplies the results in a specified order in
$U(\mathfrak g)$. Brackets are commutators. Operations on
$\operatorname{CE}(\mathfrak g\otimes E_W)$ include signed
symmetrization of suspended inputs. Two formulas are identified if
they agree for every Lie algebra with such a derivation action.
This restriction defines a universal algebra of operations; it does
not define all maps of geometric coefficient families.

Put $P=W^*$. Denote the natural convolution algebra by
$\mathcal N(P;\mathfrak d)$, with differential twisted by the scalar
cochain $\tau$. Thus
\[
 \delta f=-(-1)^{|f|}fQ+\tau*f-(-1)^{|f|}f*\tau.
\]
Complete in actor weight and allow only finite sums of differential
operations at each weight and cochain degree. The ideal of positive
actor weight is $\mathcal N_+$. Natural actions and homotopies mean
its Maurer--Cartan space, using polynomial differential forms on
simplices as parameters.

For commuting derivations $D_1,\ldots,D_r$, the unary cochain
$x_{S,a}$ extracts an exterior coefficient $e_S$ and applies
$D^a/a!$. It has degree $1-|S|$. Write $t=x_{\varnothing,0}$.
If $\epsilon(B,C)$ is the exterior shuffle sign, the unary
differential is
\begin{equation}\label{eq:vol2-natural-differential}
\begin{aligned}
 \delta x_{S,a}
 ={}&-\sum_{\substack{B\sqcup C=S\\b+c=a}}
       (-1)^{|B|(|C|+1)}\epsilon(B,C)x_{B,b}x_{C,c}\\
    &+tx_{S,a}-(-1)^{1-|S|}x_{S,a}t.
\end{aligned}
\end{equation}
It follows from
$Q_2(s(e_Bx)s(e_Cy))=(-1)^{|B|}s(e_Be_C[x,y])$
and divided Leibniz. The two terms containing $t$ cancel when
$x_{S,a}\ne t$, while $\delta t=t^2$.
For one actor and one derivation,
\[
 \delta x_{\{1\},1}=x_{\{1\},0}x_{\varnothing,1}
                      -x_{\varnothing,1}x_{\{1\},0},
 \qquad \delta x_{\varnothing,2}=-x_{\varnothing,1}^2.
\]
These identities already show the polynomial actor coordinate and
the exterior derivation class in cohomology.

\begin{theorem}\label{thm:vol2-natural-cohomology}
There is a surjective filtered dg algebra map
\begin{equation}\label{eq:vol2-natural-projection}
 \pi:\mathcal N(P;\mathfrak d)\longrightarrow
          \prod_{q\geq0}\operatorname{Sym}^q(P)
                     \otimes\Lambda(\mathfrak d[-1])
\end{equation}
with acyclic kernel in each actor weight. The target differential
is zero; its actor coordinates have degree zero and its derivation
generators degree one. On positive actor weights it induces an
equivalence of Maurer--Cartan spaces. Their components are
\begin{equation}\label{eq:vol2-natural-classes}
 \prod_{q\geq1}\operatorname{Sym}^q(W^*)\otimes\mathfrak d.
\end{equation}
At each point the fundamental group is the additive group
$\mathfrak m_W=\prod_{q\geq1}\operatorname{Sym}^q(W^*)$,
and all higher homotopy groups vanish.
\end{theorem}

\begin{proof}
Unary operations freely generate the convolution algebra. To prove
this, evaluate on the free Lie algebra on named copies of the free
$\mathfrak d$-module $U(\mathfrak d)$. In its tensor envelope,
an ordered monomial records the input, its exterior label and the
operator applied to it. Distinct such words are independent.
Polarization recovers repeated exterior types with the suspended
signs. A cochain of actor weight $q$ and degree $p$ uses $q+p$
inputs.

The Hopf coproduct of $U(\mathfrak d)$ expresses its Leibniz rule.
Rescaling the exterior-label generator by
$(-1)^{\binom{|S|}{2}}$ identifies the differential with the cobar
differential of $\Lambda(P[1])\otimes U(\mathfrak d)$, apart from
the scalar generator $t$. Thus the natural algebra is
\begin{equation}\label{eq:vol2-natural-cobar}
 (k\langle t\rangle,\delta t=t^2)
       *\Omega\bigl(\Lambda(P[1])\otimes U(\mathfrak d)\bigr).
\end{equation}
Here $*$ is free product, with the stated weight completion.
The augmentation ideal of the first factor contracts by
$h(t^{2j})=t^{2j-1}$ and $h(t^{2j-1})=0$. In the free product,
contract the first such factor of each alternating word, with
the sign of its prefix. This removes $t$ without changing
cohomology.

PBW symmetrization identifies $U(\mathfrak d)$ with the symmetric
coalgebra on $\mathfrak d$. This is a coalgebra identification,
not an identification of products. The cobar projection sends
primitive actor and derivation generators to their namesakes and
all other generators to zero. The cobar differential maps to
graded commutators or odd squares, so the projection is a dg
algebra map.

Its cohomology is computed in finite auxiliary weights. The dual
algebra of $\Lambda(P[1])\otimes\operatorname{Sym}^c(\mathfrak d)$
is exterior on degree-one variables $\alpha_i$ and polynomial on
degree-zero variables $z_j$. A free resolution of its scalar
module is
\[
 \Lambda(\alpha_i)\otimes k[z_j]
       \otimes k[v_i]\otimes\Lambda(\eta_j),\qquad
 |v_i|=0,\quad |\eta_j|=-1,\quad
 dv_i=\alpha_i,\quad d\eta_j=z_j.
\]
For an odd pair, $d(v^{a+1})=(a+1)\alpha v^a$; for an even
pair this is the two-term Koszul resolution. Each pair has scalar
cohomology, so their tensor product resolves $k$. Applying Hom
into $k$ gives zero differential and the algebra on the right
of \eqref{eq:vol2-natural-projection}. The bar resolution computes
the same Ext. The primitive projection is surjective on
cohomology and both sides have the same finite dimension in each
auxiliary weight, proving the quasi-isomorphism. For an
infinite-dimensional derivation algebra take the filtered union
of finite-dimensional symmetric-coalgebra computations.

To lift a Maurer--Cartan element, solve successively in actor
weight. At the first unsolved weight, its curvature is a closed
degree-two element of the kernel of $\pi$. A degree-one primitive
removes it. The same argument with forms relative to a simplex
boundary lifts every compatible simplex. Tensoring a weightwise
acyclic complex with these forms preserves acyclicity over $k$.
The resulting fibres are contractible. In the target, the
commutator bracket and differential vanish. Its degree-one
elements give \eqref{eq:vol2-natural-classes}, its degree-zero
elements give the stated loops, and its negative degrees vanish.
\end{proof}

For one derivation, retain also the derivative order. The sector
where actor weight and derivative order both equal $q$ has
cohomology
\[
 \operatorname{Sym}^q(P)\otimes\Lambda^q(kD)
 \quad\text{in degree }q.
\]
It is therefore acyclic for $q\geq2$. For several commuting
derivations, the same formula has their span in place of $kD$;
its degree-two term need not vanish.

\begin{lemma}\label{lem:vol2-natural-bar-lifting}
Restriction along
$\operatorname{CE}(\mathfrak h)\to\operatorname B U(\mathfrak h)$
is a surjective quasi-isomorphism on natural convolution complexes
twisted by the scalar action. Its kernel contracts naturally,
preserving actor weights and derivative orders. It has the same
property relative to any union of actor unit faces.
\end{lemma}

\begin{proof}
For abelian $\mathfrak h$, multilinear PBW weight $d$ of the bar
is the finite complex of ordered partitions of $d$ named inputs,
with differential merging adjacent blocks. The Koszul resolution
used above identifies its homology with suspended symmetrization.
Choose rational complements to cycles and boundaries and average
over $S_d$. Inverting the differential on the boundary complement
gives a permutation-equivariant contraction, normalized to
\[
 pj=1,\quad hj=ph=h^2=0,\qquad dh+hd=1-jp.
\]
Evaluation on tensor powers gives a natural contraction for every
graded $\mathfrak h$, compatible with its internal differential.

For nonabelian $\mathfrak h$, PBW identifies the change in the
bar differential with a perturbation $b$ that lowers primitive
degree. The formulas
$p'=p(1+bh)^{-1}$ and $h'=h(1+bh)^{-1}$ are finite on each
input. The inclusion stays $j$, since $bj=jQ_2$ and $hj=0$.
The transferred differential is the Chevalley differential.
Precomposition by $p'$ splits restriction, and
$f\mapsto(-1)^{|f|}fh'$ contracts its untwisted kernel.

The scalar twist preserves that kernel because $j$ is a
coalgebra map. Its commutator raises the least primitive input
degree on which a cochain is nonzero. Perturbing the kernel
contraction again therefore gives finite sums on each input.
All these maps use permutations, PBW operations, products and
brackets. They commute with exterior face projections and
preserve derivative orders. Restriction to cochains vanishing
on the chosen faces proves the relative assertion.
\end{proof}

Natural identities evaluate on the coefficient category because
they are identities of these multilinear operations. On a
configuration of $m$ positions with $n$ exterior generators,
actor weight is at most $mn$. Their formal weight products
introduce no completion of a state or of a normal derivative.

\begin{theorem}\label{thm:vol2-derivation-multiplicative}
For a vertex derivation $D$ commuting with translation, the
action $\alpha_D$ admits a coherent unital multiplicative
structure for the ordinary exterior product of $C_E$.
Within natural operations of equal actor and derivative weight,
the higher fillings are contractible after fixing the unary
action and a compatible binary bar path.
Every profile $\alpha_{\mathbf t}$ has a map of pointed actions
to this multiplicative action. In particular its two unit
insertions have the common recipient
\begin{equation}\label{eq:vol2-profile-recipient}
 \mu(\Sigma_{\mathbf t}\vot\Sigma_{\mathbf t}):
 (C_E,\alpha_{\mathbf t})^{\circledast2}
                         \longrightarrow(C_E,\alpha_D).
\end{equation}
\end{theorem}

\begin{proof}
Use the natural bar complex for $n$ exterior actors. Restriction
to their proper unit faces is surjective: the coordinate
projections commute, so inclusion--exclusion extends compatible
face data. Its kernel is the fully mixed part. For $n\geq2$
that part has actor weight at least two and is acyclic in the
equal derivative-weight sector, by
Theorem~\ref{thm:vol2-natural-cohomology} and
Lemma~\ref{lem:vol2-natural-bar-lifting}.

Consequently compatible Maurer--Cartan data on the actor faces
and the boundary of a parameter cube extend across the cube.
Lift the data linearly, then kill the first remaining curvature
in each actor weight by a primitive in the fully mixed complex
tensored with the forms relative to the cube boundary. The
Bianchi identity makes the curvature closed. Further parameter
simplices give the same lifting, proving contractibility of
the relative choices.

The component for $n$ actors uses the cube $[0,1]^{n-1}$.
At $t_i=0$ multiply the adjacent actors; at $t_i=1$ compose
the two shorter bar maps. An end unit deletes its adjacent
parameter. An interior unit combines the two parameters by
\[
 a\vee b=a+b-ab.
\]
This polynomial operation is associative, with identity zero
and absorbing element one. It makes all repeated unit
restrictions agree with the endpoint faces. The preceding
lifting constructs the binary path and then every higher
component, or extends a specified compatible binary path.

The endpoint compositions are the geometric ones. Indeed,
$\vid+hD$ over $k[h]/(h^2)$ is a vertex automorphism, and
Lemma~\ref{lem:vol2-strong-boundary-tensor} carries $D$ to
$D\otimes1$ on exterior coefficient extensions. Its naturality
identifies every bar operation and contraction with the
corresponding coefficient map. Thus the constructed families
retain collisions, normal transfers and the vacuum.

Proposition~\ref{prop:vol2-profile-actor} supplies each incoming
map. Composing its tensor square with $\mu$ gives
\eqref{eq:vol2-profile-recipient}. On either unit axis it is
$\Sigma_{\mathbf t}$ with its action path. The unital coherence
identifies these restrictions on their common vacuum, supplying
the required relative two-simplex. The same higher product
coherences receive any finite ordered product of profile maps.
\end{proof}

The contractibility is in the specified natural-operation
complex. The space of all geometric multiplicative structures
can have additional components or homotopy groups.

\section{First-weight classes and actor size}
\label{sec:vol2-first-weight}

Fix $\kappa\in k^\times$ and the Heisenberg vertex algebra
\[
 V_\kappa=k[x_1,x_2,\ldots],\qquad
 \mathsf j=x_1,\qquad x_n=\mathsf j_{-n}\mathbf1.
\]
Its states are polynomials. Translation and the positive current
modes are
\[
 Tx_n=nx_{n+1},\qquad
 \mathsf j_n=\kappa n\partial_{x_n}\quad(n>0).
\]
Let $L_0x_n=nx_n$, extended as a derivation of the polynomial
algebra. Put $\mathfrak g=L_{V_\kappa}$, $A=\vU\mathfrak g$,
$P_w=\vSym^w\mathfrak g$ and $C=C_{V_\kappa}$. Lie-module
actions on $P_w$ and $A$ are adjoint actions. A coefficient
derivation is a degree-zero map of the full family $\mathfrak g$
which obeys Leibniz for its Lie bracket. Their Lie algebra is
$\mathfrak d_V=\operatorname{Der}_{\mathscr C}(\mathfrak g)$.

In left-module coordinates, the singleton coefficient has
connection $\partial_y+T$. A coefficient derivation is therefore
represented there by a $k[y]$-linear map $B(y)$ with
\begin{equation}\label{eq:vol2-horizontal-derivation}
 \partial_yB+[T,B]=0,
\end{equation}
together with compatibility with the chiral bracket. Each input
state has a finite polynomial image. A common polynomial-degree
bound on all inputs is not imposed. For example,
$B(y)=L_0+yT$ is horizontal because $[T,L_0]=-T$. The conformal
identity for $Y(a,z)$ gives
\[
 B(q)Y(a,p-q)b
   =Y(B(p)a,p-q)b+Y(a,p-q)B(q)b.
\]
Taking arbitrary Laurent principal parts proves its coefficient
derivation identity, including the normal derivatives.

The current shift $D=\partial_{x_1}$ is another such derivation:
it is the derivative of $\mathsf j\mapsto\mathsf j+t\mathbf1$.
Every state zero mode $a_{(0)}$ is a vertex derivation commuting
with $T$, by the mode commutator identity. Put
$Z_r=(\mathsf j^r)_{(0)}$ for $r\geq2$. Normal ordering gives
\begin{equation}\label{eq:vol2-Heisenberg-derivations}
 Z_2=2\kappa T,\qquad
 Z_r\mathsf j=r(r-1)\kappa x_1^{r-2}x_2,
 \qquad (\mathsf j^2)_{(1)}=2\kappa L_0.
\end{equation}
Indeed the contraction of $\mathsf j^r(z)$ with one current is
$r\kappa z^{-2}{:}\mathsf j(z)^{r-1}{:}$, whose residue is
$r\kappa T(\mathsf j^{r-1})$. There is no multiple contraction
with the single current on the right. The quadratic zero and
first modes give the displayed formulas for $T$ and $L_0$ on
every polynomial state. The images of $\mathsf j$ show that
$D,Z_2,Z_3,\ldots$ are linearly independent.

\subsection{An additional output root}

The family $P_w$ is the open direct image of its coefficient
on $w$ distinct roots, supported on configurations with at most
$w$ distinct positions. On the complement of the locus with
fewer roots, each ordered partial diagonal is a copy of
$\operatorname{Conf}_w(X)$. Its coefficient is the tensor of
$w$ state modules, equivariant under $S_w$. Reconstruction is
exact: use closed direct image on each partial diagonal and
localization along its cross-root divisors. This also retains
every normal derivative.

\begin{lemma}\label{lem:vol2-root-support-bound}
For $a,w\geq1$,
\[
 \operatorname{Ext}^q_{\mathfrak g}(P_w,P_a)=0
                       \quad\text{if }q<a-w.
\]
For every coefficient derivation $B$, the formula
\begin{equation}\label{eq:vol2-extra-root}
 \gamma_B^{(w)}(x;v_1\cdots v_w)
        =m_{w+1}(Bx,v_1,\ldots,v_w)
\end{equation}
defines a class in $\operatorname{Ext}^1_{\mathfrak g}(P_w,P_{w+1})$,
where $m_{w+1}$ is symmetric multiplication. The map from
$\mathfrak d_V$ to this extension group is injective.
\end{lemma}

\begin{proof}
A coefficient family supported on at most $a-1$ roots has zero
derived Hom to $P_a$, by open adjunction. The $n$th Chevalley
term computing the desired module maps is
\[
 \vRH_{\mathscr C}(\Lambda^n_\star\mathfrak g\star P_w,P_a)[-n].
\]
It vanishes for $w+n<a$. Before the shift it has no negative
cohomology, since its inputs are coefficient modules in degree
zero. The first-quadrant totalization proves the bound.

Give $P_{w+1}\oplus P_w$ the triangular action with
off-diagonal term \eqref{eq:vol2-extra-root}. Its module defect
is symmetric multiplication by
$[x,By]+[Bx,y]-B[x,y]$, hence is zero on the full coefficient
diagram. Open adjunction also gives
$\operatorname{Hom}_{\mathscr C}(P_w,P_{w+1})=0$.
The underlying split extension has only its canonical
coefficient splitting. If $B\ne0$, this splitting is not
equivariant: on $w+1$ distinct roots put a state on which $B$
is nonzero at the acting root and the vacuum at all other
roots. Its image under \eqref{eq:vol2-extra-root} is nonzero.
The same argument for differences proves injectivity.
\end{proof}

To compare homotopies, resolve the pointed arrow before taking
relative convolution cochains. Bar--cobar duality gives maps
under $A$ from $U_\star(\mathfrak g\ltimes e\mathfrak g)$ to
$A$. At the scalar action, the weight-$w$ complex is
\begin{equation}\label{eq:vol2-action-tangent}
 T_w=\vRH_{\mathfrak g}(P_w,A).
\end{equation}
Indeed the underlying relative Chevalley comodule is
$\vCE(\mathfrak g)\star\operatorname{Sym}^c_
 \star(e\mathfrak g[1])$, and $e\mathfrak g[1]=\mathfrak g$.
The coefficient resolutions are taken in the full enhanced
diagram. This computes the homotopy fibre specifying the
pointing, rather than the kernel of an unresolved restriction.

Scaling $e$ by a parameter over $k[t]/(t^2)$ defines
\[
 \nu:\pi_0\Phi_C(C_E)\longrightarrow H^1(T_1).
\]
The first coefficient is a cocycle; a pointed homotopy changes
it by a boundary. More generally a relative homotopy has a
normalized convolution gauge $G=1+g_1+g_2+\cdots$ satisfying
$d_0G=G\theta-\theta'G$. This follows by integrating the
polynomial interval equation successively in positive weight.
Every inverse and exponential is finite on a fixed
configuration. If both actions vanish in weights below $r$,
the equations give
\begin{equation}\label{eq:vol2-leading-gauge}
 d_0g_j=0\quad(j<r),\qquad
 d_0g_r=\theta_r-\theta'_r.
\end{equation}
The lower gauge coefficients need not be zero.

\begin{theorem}\label{thm:vol2-first-weight-classification}
There are vector-space identifications
\begin{equation}\label{eq:vol2-first-weight-classification}
 \operatorname{Ext}^1_{\mathfrak g}(P_1,P_2)
       \cong\mathfrak d_V,
 \qquad H^1(T_1)=k\xi\oplus\mathfrak d_V.
\end{equation}
The first identification is $B\mapsto[\gamma_B^{(1)}]$.
The class $\xi$ changes the zero-mode charge at one root.
For every $B\in\mathfrak d_V$, the tree action $\alpha_B$
has first class $[\gamma_B^{(1)}]$.
\end{theorem}

\begin{proof}
In the Chevalley extension spectral sequence
\[
 E_1^{n,q}=\operatorname{Ext}^q_{\mathscr C}
      (\Lambda^n_\star\mathfrak g\star\mathfrak g,P_2),
 \qquad n,q\geq0,
\]
open adjunction kills the $n=0$ column. Only $(n,q)=(1,0)$
can contribute in total degree one. There are no incoming
boundaries, and every higher differential has a negative index.
Thus the extension group consists precisely of module cocycles
$\Gamma:\mathfrak g\star\mathfrak g\to P_2$.

On the two-root locus write $R=k[p,q,(p-q)^{-1}]$. Such a
map is an $R$-linear endomorphism of
$R\otimes(V_\kappa\otimes V_\kappa)$ with
\[
 \partial_p\Gamma+[T\otimes1,\Gamma]=0,
 \qquad\partial_q\Gamma+[1\otimes T,\Gamma]=0.
\]
Its individual input images are finite; no uniform pole bound
is assumed. In the module cocycle equation give the two acting
inputs positions $r,p$ and the module input position $q$.
Away from the triple diagonal the target has three disjoint
pieces $r=q$, $p=q$, $r=p$, whose open direct images detect
equality.

On $r=q$, comparison of the zeroth normal coefficients after
testing $(r-q)^n$ gives
\[
 \Gamma(y\otimes x_{(n)}v)=(1\otimes x_{(n)})\Gamma(y\otimes v)
 \quad(n\geq0).
\]
The full normal transfer is retained; this projection supplies
necessary commutation identities. Take $x=\mathsf j$, $n>0$,
and $x=\mathsf j^2/(2\kappa)$, $n=1$.
Then $\Gamma$ commutes with all partial derivatives and $L_0$
in the second state factor. Its restriction to the weight-zero
line of that factor defines $B$. By induction on second-state
weight, $\Gamma(a\otimes v)-Ba\otimes v$ is killed by all
second-factor partial derivatives. It is therefore constant
in that factor, while having positive weight, and is zero.
Hence $\Gamma=B(p,q)\otimes1$.

The second horizontality equation gives $\partial_qB=0$.
A $q$-invariant element of $R$ lies in $k[p]$: it depends only
on $p$ in the fraction field, and substitution $q=p-1$ makes
it polynomial. This applies to each matrix coefficient, so
$B=B(p)$ and it obeys \eqref{eq:vol2-horizontal-derivation}.
On $r=p$ the cocycle equation is the derivation identity for
$B$, tensored with the unchanged spectator at $q$. The normal
transfer differentiates its coordinate dependence; $B(r)$ is
not replaced by $B(p)$ before that transfer. Tensoring with a
spectator vacuum and inverting $q-p$ is faithful, so this
identity holds in the full coefficient family. This proves
surjectivity; Lemma~\ref{lem:vol2-root-support-bound} proves
injectivity.

Equivariant PBW decomposes $A$ by output number. It is both a
sum and a product, because only finitely many output numbers
occur at any configuration. Lemma~\ref{lem:vol2-root-support-bound}
leaves only one and two outputs in $H^1(T_1)$. At one root,
extensions are extensions of the unital moving Heisenberg
vacuum. The positive modes give the polynomial Koszul
contraction for $\kappa n\partial_{x_n}$. Its nonconstant
monomials contract by the Euler homotopy, leaving the zero-mode
class. The de Rham complex of the affine root coordinate has
only its scalar class, so this contribution is $k\xi$.
Concretely, over $k[\epsilon]/(\epsilon^2)$ take
$\mathsf j_0=\epsilon$ and connection
$\partial_y+T+\epsilon x_1/\kappa$. Its commutator with
$\mathsf j_1$ is $-\epsilon$, and the other current translation
relations remain unchanged. This is the charge extension;
its zero-mode action prevents a split. These calculations
identify the one-output summand.

Finally the first tree term is $Bx\,v$. PBW writes it as
$m_2(Bx,v)+[Bx,v]/2$. For the zero-cochain $B$,
$(dB)(x;v)=[x,Bv]-B[x,v]=-[Bx,v]$.
The output-one term is exact and the first class is the
asserted output-two extension.
\end{proof}

Horizontality in the proof does not mean that an arbitrary
two-root map has bounded poles on the entire state space.
For $z=p-q$, the operator
\begin{equation}\label{eq:vol2-horizontal-poles}
 \mathcal K_z=\kappa\sum_{m,n\geq1}
  \frac{(-1)^{m-1}(m+n-1)!}{(m-1)!(n-1)!}
      z^{-m-n}\partial_{x_m}\otimes\partial_{x_n}
\end{equation}
is horizontal and finite on each polynomial input. For its
scalar coefficients $c_{m,n}$,
$m c_{m+1,n}=-(m+n)c_{m,n}$ and
$n c_{m,n+1}=(m+n)c_{m,n}$; together with
$[T,\partial_{x_m}]=-(m-1)\partial_{x_{m-1}}$ these prove
the two connection equations. Its pole order on
$x_m\otimes x_n$ grows with $m+n$. The commutant condition,
not a uniform pole assumption, is what removes the poles in
Theorem~\ref{thm:vol2-first-weight-classification}.

\subsection{Bounds on an actor}

\begin{corollary}\label{cor:vol2-unbounded-arity-actions}
For every $N\geq1$ there is a nontrivial pointed Heisenberg
action which is the trivial action on every configuration with
at most $N$ positions, including all its collisions and normal
derivatives.
\end{corollary}
\begin{proof}
Take $t_N=1$ and the other profile parameters zero. Its first
nonzero cochain is $Dx\,u_{[N]}$. Its PBW output-$N+1$
component is the nonzero extension
$[\gamma_D^{(N)}]$ of
Lemma~\ref{lem:vol2-root-support-bound}.
Equation~\eqref{eq:vol2-leading-gauge} prevents a homotopy to
the trivial action. Every nonconstant term uses at least
$N+1$ positive supports, so vanishes on the stated smaller
configurations.
\end{proof}

\begin{proposition}\label{prop:vol2-higher-weight-product}
Suppose a pointed action of $C_E$ has zero positive weights
below $r>1$ and nonzero first class $[\theta_r]$. It cannot
be multiplicative for the standard exterior-algebra product.
\end{proposition}
\begin{proof}
Pull the two candidate binary actions back along the primitive
diagonal induced by $e\mapsto e_1+e_2$. Multiplication followed
by the action has first class $2^r[\theta_r]$. Composition has
first class $2[\theta_r]$. To see the latter with its paths,
scale the two actors independently by $s,t$ and truncate at
total degree $r$. The first nonzero jet of composition is
the sum of the pure jets, since a product of two perturbations
has degree at least $2r$. Unit homotopies identify the pure
terms. A pointed binary path would therefore force
$(2^r-2)[\theta_r]=0$ by
\eqref{eq:vol2-leading-gauge}, contrary to characteristic zero.
\end{proof}

The profile actions satisfy this obstruction when their first
parameter has index greater than one. They still have the
multiplicative recipient \eqref{eq:vol2-profile-recipient}.

For any pointed factorizing actor $F$, let $B_F$ be a cobar
model under $U_\star L_k$, with its induced adjoint $L_k$-action.
At the augmented map $C_E\to J\to F$, the first-weight
mapping complex is
\[
 T_1^{E,F}=\vRH_{L_k}(L_k,B_F).
\]
The relative Chevalley argument of
\eqref{eq:vol2-action-tangent} gives this identification,
without a degree bound on $B_F$. A fixed action $\beta$ of
$F$ induces a linear map
\begin{equation}\label{eq:vol2-action-chain-rule}
 b_\beta:H^1(T_1^{E,F})\longrightarrow H^1(T_1),\qquad
 \nu\bigl(\beta(f\vot\vid_C)\bigr)=b_\beta(\nu_F(f)).
\end{equation}
Indeed, terms with one perturbation in a compatible bar
composite give a chain map; terms with two perturbations
vanish over $k[t]/(t^2)$. Resolving the pointed arrow retains
its prescribed homotopy. This finite free scalar extension
commutes with derived coefficient maps and diagram limits.

\begin{lemma}\label{lem:vol2-finite-exterior-tangent}
For $W$ finite-dimensional in degree zero and
$F=C_{E_W}$, one has $H^1(T_1^{E,F})=W$.
Its vectors are represented by the linear exterior maps
$e\mapsto w$.
\end{lemma}
\begin{proof}
In PBW for $U_\star L_{E_W}$, let $a$ count output roots
and $b$ their total exterior degree. The support bound
puts Chevalley degree at least $a-1$. Total degree one
allows only $(a,b)=(1,0),(1,1),(2,0)$. The first is
$H^1_{\mathrm{dR}}(\mathbb A^1;k)=0$. The second is
$\operatorname{Hom}_{L_k}(L_k,L_k)\otimes W=W$.
For $(2,0)$, open adjunction leaves only the scalar
multiple $m$ of symmetric multiplication in Chevalley
degree one. Its differential is
$(dm)(x,y;v)=m([x,y],v)$, nonzero on a simple pole between
the acting points with the module root distinct. The vacuum
principal-part bracket has residue one there. Thus this
last summand supplies no class.
\end{proof}

\begin{theorem}\label{thm:vol2-finite-actor-obstruction}
If a pointed action $(F,\beta)$ represents $\Phi_C$, then
$H^1(T_1^{E,F})$ is infinite-dimensional. The same conclusion
holds if $(F,\beta)$ is a local centre in $\vMon$.
Consequently no coalgebra equivalent as a pointed object
to $C_{E_W}$, with $W$ finite-dimensional and $d=T=0$
on its exterior generators, can underlie either object.
\end{theorem}
\begin{proof}
Representability supplies an incoming map from every
$\alpha_{Z_r}$. Equation~\eqref{eq:vol2-action-chain-rule}
puts every $[\gamma_{Z_r}^{(1)}]$ in the image of $b_\beta$.
These classes are linearly independent by
\eqref{eq:vol2-Heisenberg-derivations} and
Theorem~\ref{thm:vol2-first-weight-classification}.
If only finality in $\vMon$ is assumed, the same argument
applies because each $\alpha_{Z_r}$ is multiplicative by
Theorem~\ref{thm:vol2-derivation-multiplicative}.
Lemma~\ref{lem:vol2-finite-exterior-tangent} excludes the
stated finite exterior actors in both cases.
\end{proof}

A finite exterior actor can still receive a specified finite
family. The theorem excludes its universal property for all
actions. It does not exclude a larger factorizing centre or
identify the whole formal deformation functor of a recipient.

\section{Recipients for nilpotent families}
\label{sec:vol2-nilpotent-recipients}

Several derivations need not commute. Their action classes compose
by the Baker--Campbell--Hausdorff law. An exterior actor can carry
that law, but its multiplication must retain which input supplies
each output generator. We first describe these marked coefficients.

\subsection{Natural actor maps}

In this subsection all actor maps are under and over $J$, with both
comparison paths. Let $W,V$ be finite-dimensional vector spaces
in degree zero, and write $C_W=C_{E_W}$ and $C_V=C_{E_V}$.
Natural operations between $C_W$
and $C_V$ use products, brackets, permutations and linear maps of
exterior coefficients. Each ordinary Lie input is used once.
No additional derivation acts on it. The relative convolution
cochains vanish on purely ordinary inputs and take values in the
augmentation ideal of $U_\star L_{E_V}\to U_\star L_k$.

\begin{theorem}\label{thm:vol2-natural-actor-maps}
The space of natural maps $C_W\to C_V$ is equivalent to the
discrete space of formal matrices
\[
 A(t)\in\operatorname{Hom}(W,V)[[W^*]].
\]
The coefficient of polynomial degree $r-1$ has source actor
weight $r$. Its coordinate contraction is $\Phi_A(t)=A(t)t$.
Composition and diagonal tensor are
\begin{equation}\label{eq:vol2-marked-composition}
 (B\circ A)(t)=B(A(t)t)A(t),\qquad
 (A\oplus B)(t,s)=
 \begin{pmatrix}A(t)&0\\0&B(s)\end{pmatrix}.
\end{equation}
The identity is the constant identity matrix.
\end{theorem}

\begin{proof}
Let $K_W=(E_W^+)^\vee$ and let $\mathfrak q_W$ be its Lie
cobar. For a basis of $W$, a generator $y_I$ indexed by a
nonempty exterior subset has degree $1-|I|$ and differential
\[
 dy_I=\frac12\sum_{\substack{B\sqcup D=I\\B,D\ne\varnothing}}
       (-1)^{|B|-1}\epsilon(B,D)[y_B,y_D].
\]
The sign is obtained by the exterior-label rescaling used in
\eqref{eq:vol2-natural-cobar}. In particular
$dy_{\{i,j\}}=[y_i,y_j]$ for $i<j$.
Adjoin a degree-one Lie generator $z$ with
$dz=[z,z]/2$ by Lie coproduct, writing
$\mathfrak h_W=\operatorname{Lie}(z)\amalg\mathfrak q_W$.

The natural convolution algebra identifies with the summand
of $\widehat{U(\mathfrak h_W\otimes E_V)}$ having positive
source weight and positive target exterior degree. To check
this assertion, polarize the inputs and evaluate on the free
Lie algebra on their names. A Lie word specifies an output
root, and a PBW product specifies a forest of roots, each with
its exterior label. Distinct forests give independent
multilinear operations. Evaluation on the group-like exponential
of the suspended inputs turns convolution into multiplication.
Transposing exterior multiplication gives the displayed Lie
cobar differential. The scalar twist cancels the couplings
between $z$ and the nonempty labels by equivariance; its
remaining equation is $dz=z^2$.

Killing $z$ is a quasi-isomorphism in each source weight.
Before tensoring with $E_V$, the associative envelope is
\[
 (k\langle z\rangle,dz=z^2)*U(\mathfrak q_W).
\]
Contract the first positive $z$ factor in every alternating
word by $h(z^{2j})=z^{2j-1}$, $h(z^{2j-1})=0$, with the
prefix sign. PBW identifies the Lie algebra with a direct
summand of its envelope, so the Lie projection is also a
quasi-isomorphism. Tensoring with $E_V$ and taking envelopes
preserve it: tensor powers and symmetric coinvariants are
exact over $k$.

The projection $\mathfrak q_W\to W^*$ sends $y_i$ to its
namesake in an abelian Lie algebra and kills longer exterior
labels. The resolution
$E_W\otimes k[v_i]$, $dv_i=e_i$, contracts in positive
combined weight $N$ by $N^{-1}\sum_i v_i\partial_{e_i}$.
Dualizing the bar comparison in each finite weight proves
$H(U\mathfrak q_W)=\operatorname{Sym}(W^*)$, with this
projection. The PBW length-one summand again gives the Lie
quasi-isomorphism. Consequently the relative natural algebra
has a surjective dg algebra projection, with weightwise
acyclic kernel, to
\[
 \mathcal B(W,V)=
 \left[\widehat{\operatorname{Sym}}(W^*\otimes E_V)
                         \right]_{E_V^+}.
\]
The subscript denotes the ideal generated by the positive
target exterior labels. Its differential is zero. It has
no terms in degrees at most zero and
\[
 \mathcal B^1(W,V)=
 \widehat{\operatorname{Sym}}(W^*)\otimes W^*\otimes V.
\]
Relative Maurer--Cartan lifting, as in
Theorem~\ref{thm:vol2-natural-cohomology}, has contractible
fibres. In the target a degree-one simplex is a polynomial
function with zero differential, hence constant. This proves
the claimed discreteness.

The surviving $W^*\otimes V$ factor marks the input supplying
the target exterior generator. The other input factors are
unmarked. Exponential polarization supplies their factorials.
In a bar composite, unmarked coefficients substitute
$t\mapsto A(t)t$, while the marked factor is transported by
$A(t)$ and then by $B(A(t)t)$. All joined-root terms vanish
under the projection. This proves the first formula in
\eqref{eq:vol2-marked-composition}; disjoint actor groups give
the block-diagonal formula. The relative bar contraction
transports these calculations to actual maps and paths.
At $m$ positions the source actor weight is at most
$m\dim W$, so every formula is finite there.
\end{proof}

The theorem classifies the specified natural maps. Full
geometric mapping spaces can contain coefficient extensions.
For example, with left differential modules and
$D_x=k\langle x,\partial_x\rangle$, put
$M=k[x,x^{-1}]$ and $\Delta=D_x/D_xx$. The presentation
$M=D_x/D_x(x\partial_x+1)$ gives a two-term free resolution.
On $\partial_x^n\delta\in\Delta$, the resulting Hom
differential $x\partial_x+1$ is multiplication by $-n$.
Thus $\operatorname{Hom}_{D_x}(M,\Delta)=k$ and
$\operatorname{Ext}^1_{D_x}(M,\Delta)=k$. This normal-chart
extension alone does not construct a compatible operation on
the whole collision diagram.

\subsection{Multiplication with a chosen frame}

An associative multiplication on $C_W$ has a matrix
$M(x,y)=(L(x,y)\mid R(x,y))$ and coordinate contraction
$F(x,y)=L(x,y)x+R(x,y)y$. Its unit conditions are
$L(x,0)=I$ and $R(0,y)=I$. By
\eqref{eq:vol2-marked-composition}, associativity is precisely
\begin{equation}\label{eq:vol2-multiplication-blocks}
\begin{aligned}
 L(F(x,y),z)L(x,y)&=L(x,F(y,z)),\\
 L(F(x,y),z)R(x,y)&=R(x,F(y,z))L(y,z),\\
 R(F(x,y),z)&=R(x,F(y,z))R(y,z).
\end{aligned}
\end{equation}
It gives the fibrewise-linear formal group law
\[
 (x,u)(y,v)=(F(x,y),L(x,y)u+R(x,y)v).
\]
Both its zero section and its diagonal are homomorphisms.
Conversely, such a law gives the displayed matrix equations.
Discreteness of the natural mapping spaces supplies a
contractible space of higher monoid coherences.

Let $\mathfrak a$ be the tangent Lie algebra of $F$.
A compatible left-symmetric product on $\mathfrak a$ is
$x\cdot y=\ell_xy$ such that
\[
 \ell_xy-\ell_yx=[x,y],\qquad
 [\ell_x,\ell_y]=\ell_{[x,y]}.
\]
These products describe the fibre laws up to changes of frame
which fix the diagonal and are the identity on the fibre at zero.
Here is the construction, including the frame itself.

The affine Lie map
\[
 x\longmapsto\begin{pmatrix}\ell_x&x\\0&0\end{pmatrix}
\]
integrates formally to a representation $\rho_0$ and a cocycle
$b$, with $db_0=I$ and
$b(F(x,y))=b(x)+\rho_0(x)b(y)$.
In exponential coordinates these are
$\rho_0(x)=e^{\ell_x}$ and
$b(x)=\int_0^1e^{t\ell_x}x\,dt$. The formal matrix
exponential and its Baker--Campbell--Hausdorff identity prove
both equations. Choose an invertible frame $M(x)$ with
$M(0)=I$ and $M(x)x=b(x)$. One always exists: if
$b=\sum_{r\geq1}b_r$ is its homogeneous expansion, take
$M=\sum_{r\geq1}Db_r/r$ and use Euler's identity.
Then
\begin{equation}\label{eq:vol2-graded-frame}
 L(x,y)=M(F(x,y))^{-1}M(x),\qquad
 R(x,y)=M(F(x,y))^{-1}\rho_0(x)M(y)
\end{equation}
satisfies \eqref{eq:vol2-multiplication-blocks} and has
coordinate contraction $F$.

For the converse, the zero section splits the fibre group.
Multiplication $(0,w)(x,0)=(x,L(0,x)w)$ trivializes it.
The frame $M(x)=L(0,x)^{-1}$ transforms the law to a
semidirect product, and the diagonal becomes a cocycle
$b(x)=M(x)x$. Differentiating its identities gives exactly
the two identities for $\ell$. Another frame $M'$ changes
the original law by $S(x)=M'(x)^{-1}M(x)$, with $S(x)x=x$
and $S(0)=I$. Conversely such an isomorphism leaves the
infinitesimal product unchanged. This proves the assertion
about changes of frame, while retaining the distinct
multiplication matrices on a fixed actor.

For instance, on $W=k^2$ with additive $F$, put
$N(t)=t(-t_2,t_1)$ and $S(t)=I+N(t)$. Then $N^2=0$,
$S^{-1}=I-N$ and $S(t)t=t$. The frame gives
$L(x,y)=S(x+y)^{-1}S(x)$ and
$R(x,y)=S(x+y)^{-1}S(y)$. Although its fibre conjugation
is trivial, its multiplication is not the constant additive
matrix: $L((1,0)^T,(0,1)^T)=
\left(\begin{smallmatrix}2&1\\1&1\end{smallmatrix}\right)$.

For a finite-dimensional positively graded Lie algebra
$\mathfrak a=\bigoplus_{j>0}\mathfrak a_j$, let
$G|_{\mathfrak a_j}=j$. The product
\begin{equation}\label{eq:vol2-positive-product}
 \ell_xy=G^{-1}[x,Gy]
\end{equation}
is compatible: the derivation identity for $G$ proves its
antisymmetrization, and
$\ell_x=G^{-1}\operatorname{ad}_xG$ proves its representation
identity. In Baker--Campbell--Hausdorff coordinates choose
$M(x)=\int_0^1e^{t\ell_x}\,dt$. Every $\ell_x$ raises
positive grading. The exponentials, $M^{-1}$ and $F$ are
therefore polynomials. Formula~\eqref{eq:vol2-graded-frame}
constructs a natural coherent monoid on $C_{\mathfrak a}$.

\subsection{Composition of action classes}

For a natural action with full twisting cochain $\tau$,
adjoin independent central degree-zero test elements $z_i$
annihilated by the derivations and put $\zeta=\sum_i e_i z_i$.
The operator
\[
 T_\tau(z)x=\tau\bigl(sx\exp(s\zeta)\bigr)
\]
has constant term $x$. It defines a series in
$\prod_{q\geq0}\operatorname{Sym}^q(W^*)\otimes U(\mathfrak d)$.
The twisting equation on two ordinary inputs gives
$T_\tau[x,y]=[T_\tau x,T_\tau y]$. Evaluate on the free
Lie algebra on the named free modules
$U(\mathfrak d)x\oplus U(\mathfrak d)y$. The map
$a\otimes b\mapsto[ax,by]$ is injective, by projection to
the corresponding ordered word in its tensor envelope.
The preceding identity therefore says
$\Delta T_\tau=T_\tau\otimes T_\tau$.

Its logarithm is primitive, hence belongs to the completed
$\mathfrak d$-valued series. PBW projection to degree one
identifies this logarithm with the action class $\phi$ of
\eqref{eq:vol2-natural-classes}: powers of the one element
$\phi$ are already symmetric powers. Thus $T_\tau=e^\phi$.
The actual tensor extension of an inner action carries
$D$ to $D\otimes1$, by
Lemma~\ref{lem:vol2-strong-boundary-tensor}. It fixes the
outer central test parameters. Composition with an outer
action of class $\phi$ and an inner action of class $\psi$
therefore has operator $e^\phi e^\psi$ and class
\begin{equation}\label{eq:vol2-action-BCH}
 \operatorname{BCH}(\phi,\psi)=\log(e^\phi e^\psi).
\end{equation}
The bracket is $[fD,gB]=fg[D,B]$ and raises actor weight.
For example $[D,Z_3]=6\kappa T$ and $T$ is central, so
\[
 \operatorname{BCH}(pD,qZ_3)=pD+qZ_3+3\kappa pqT.
\]
Reversing the action order reverses the last sign.

Pullback along a natural actor map of matrix $A(t)$ sends
an action class $\psi$ and a scalar loop $\sigma$ to
\begin{equation}\label{eq:vol2-class-substitution}
 \psi(A(t)t),\qquad \sigma(A(t)t).
\end{equation}
Indeed removing the marked generator gives $q_j\mapsto
\sum_iA_{ji}(t)t_i$. Precomposition on the natural bar
preserves products; its scalar derivation generators are
unchanged. This proves the formula on the whole cohomology
algebra. Postcomposition by a fixed outer action preserves
the scalar loop of the inner action: the outer operator
fixes its central, derivation-inert test coefficients.

For a finite-dimensional nilpotent Lie algebra $\mathfrak a$,
let $F(x,y)=\operatorname{BCH}(x,y)$. Its normalized formal
$q$-cochains with trivial scalar coefficients are formal series
$c(x_1,\ldots,x_q)$ which vanish whenever any argument is zero.
The completion is in total coordinate degree. Their differential is
\begin{equation}\label{eq:vol2-formal-group-differential}
\begin{aligned}
 (\delta c)(x_1,\ldots,x_{q+1})
 ={}&c(x_2,\ldots,x_{q+1})\\
 &+\sum_{i=1}^{q}(-1)^i
 c(x_1,\ldots,F(x_i,x_{i+1}),\ldots,x_{q+1})\\
 &+(-1)^{q+1}c(x_1,\ldots,x_q).
\end{aligned}
\end{equation}
Associativity and the unit give $\delta^2=0$. Degree zero is
$k$ with zero differential. Denote the cohomology by
$H^\bullet_{\mathrm{for}}(F;k)$. These are functions near the
formal identity, rather than cochains on the discrete group of
$k$-points.

\begin{theorem}\label{thm:vol2-graded-multiplicative-action}
Let $\mathfrak a$ be finite-dimensional and positively graded,
and let $\rho:\mathfrak a\to\mathfrak d_V$ be a Lie map.
There is a natural action of $C_{\mathfrak a}$ on the
Heisenberg boundary with class $x\mapsto\rho(x)$, compatible
with the monoid \eqref{eq:vol2-graded-frame}. Compatibility
includes every unit and associativity homotopy. Every natural
action in this class admits such a structure.
\end{theorem}

\begin{proof}
First work universally with derivation algebra $\mathfrak a$
and class $\theta(x)=x$. The natural convolution projection
constructs an action in this class. Its two binary composites
have the same class, since
\[
 \theta(F(x,y))=\operatorname{BCH}(\theta(x),\theta(y)).
\]
Choose a comparison path $h(x,y)$ relative to both unit axes.
The fully mixed summand is a direct summand of the acyclic
kernel of that projection, so relative Maurer--Cartan lifting
retains these unit restrictions.

Let $K_n$ be the associahedron of bracketings of $n$ ordered
inputs; its dimension is $n-2$. For a module the last input
is $C$ and the others are actors. The path $h$ occupies
$K_3$. The boundary of $K_4$ is now a scalar loop
$o(x,y,z)$. It vanishes on each unit face. Replacing $h$ by
a normalized scalar loop $\sigma(x,y)$ changes it by
\[
 (\delta\sigma)(x,y,z)=\sigma(y,z)-\sigma(F(x,y),z)
                +\sigma(x,F(y,z))-\sigma(x,y).
\]
These transports are \eqref{eq:vol2-class-substitution} and
the unchanged inner scalar loop. The boundary in the next
associahedron gives $\delta o=0$: each transported edge
occurs twice with opposite orientation, the square faces
express functoriality, and the pure actor face is supplied
by the monoid coherence. Hence $o$ defines a class in the
normalized formal-group cohomology $H^3_{\mathrm{for}}(F;k)$,
with trivial scalar coefficients.

This cohomology equals $H^3(\mathfrak a;k)$ equivariantly
for the grading torus. To see this, finite-order distributions
at the identity form $U(\mathfrak a)$, with PBW associated
graded. Formal functions are its full linear dual. The
normalized group cochain complex is the dual bar resolution
of the scalar module. The Chevalley resolution computes the
same derived Hom; its exactness follows from the polynomial
Koszul resolution under PBW. Antisymmetrization into the bar
dualizes to differentiation at the identity. This gives the
asserted comparison in every degree; see also
\cite{Lechner2012}*{Theorem~4.2 and Corollary~4.3}.

The obstruction class is independent of the chosen lift and
binary path. The grading torus preserves $\theta$,
\eqref{eq:vol2-positive-product}, its frame and the natural
bar comparisons, so this class is invariant. But every
weight of $\Lambda^q\mathfrak a^*$ is strictly negative
for $q>0$. Thus its invariant cohomology vanishes. Choose
a normalized $\sigma$ with $\delta\sigma=-o$ and correct
the binary path. The resulting boundary fills.

All later attaching spheres have dimension at least two.
Their action spaces have zero homotopy groups in these
dimensions by Theorem~\ref{thm:vol2-natural-cohomology}.
Successive module associahedra therefore fill, with their
specified unit faces. Positive grading has finite pieces,
so the formal corrections respect the completion. Evaluation
is finite at every configuration. Finally applying $U(\rho)$
to the universal operations preserves all their equations,
paths and units. A natural homotopy to another action in
the same class transports the constructed module structure.
\end{proof}

The binary path in this proof is chosen and corrected. The
theorem does not require it to equal an independently prescribed
path. Nor does the invariant vanishing assert that all third
cohomology vanishes. For $[a,b]=c$ with weights $1,1,2$,
the Chevalley differential is $dc^*=-a^*\wedge b^*$.
All degree-two cochains are closed, so
$[a^*\wedge b^*\wedge c^*]$ is a nonzero top class,
of grading weight $-4$.

\begin{corollary}\label{cor:vol2-nilpotent-recipient}
A finite family of natural exterior-actor actions whose classes
take values in one finite-dimensional nilpotent Lie subalgebra
of $\mathfrak d_V$ has a common coherent multiplicative
recipient on a finite-dimensional exterior actor. Their
classes may be formal series.
\end{corollary}
\begin{proof}
If the subalgebra is zero, use $J$. Otherwise choose its
generators and nilpotence class $c$. The free nilpotent Lie
algebra $\mathfrak a$ of class $c$ on these generators is
finite-dimensional, positively graded by bracket length, and
surjects onto the subalgebra. A linear section lifts each
formal action class to a series $\Phi(t)\in\mathfrak a[[W^*]]$
with zero constant term. If $\Phi=\sum_{r\geq1}\Phi_r$,
the matrix
\[
 A_\Phi(t)=\sum_{r\geq1}\frac1rD\Phi_r(t)
\]
satisfies $A_\Phi(t)t=\Phi(t)$. Theorem~\ref{thm:vol2-natural-actor-maps}
realizes it. Pullback of the action in
Theorem~\ref{thm:vol2-graded-multiplicative-action} has the
given class by \eqref{eq:vol2-class-substitution}, so equality
of natural classes supplies the required pointed action path.
\end{proof}

The incoming maps in this corollary need not preserve
multiplication. For additive one-dimensional actors, the
map $\Phi(t)=t^2$ has matrix $A(t)=t$, but
$\Phi(s+t)-\Phi(s)-\Phi(t)=2st$.

There is also a precise limitation on paths within this
natural augmented category. Fix an actor map $f:C_W\to C_V$
and an action path $p$ over it. For a nonzero scalar loop
$\sigma\in\mathfrak m_W$, the paths $p$ and $p\cdot\sigma$
define morphisms in different components. Their difference
persists after postcomposition through any natural exterior
actor. Indeed the path fibre has components forming a torsor
under $\mathfrak m_W$, and the actor mapping spaces are
discrete, so no actor loop identifies these components.
Postcomposition concatenates both paths with the same
pulled-back path. This obstruction concerns the specified
natural maps preserving augmentation. A geometric recipient
can have additional actor loops, whose images must be included
in the full question.

\section{The full density object}
\label{sec:vol2-action-density}

Fix an arbitrary factorizing boundary $C\in\vF$.
Consider all pointed factorizing actions on $C$, with their geometric
mapping spaces. Their ambient colimit is an action, but factorization
requires comparing one common actor index with independent indices
in the separated blocks. This comparison has a section. Its kernel
is the remaining obstruction to pointed representability.

\subsection{Ambient endomorphisms and a small index}

The enhanced coefficient category $\mathscr C$ is presentable.
Indeed its differential-module categories are presentable, its
extraordinary pullbacks preserve colimits, and the small limit
defining \eqref{eq:vol2-action-category} is a limit of presentable
categories and colimit-preserving functors. Colimits are computed
on the coefficients. The same assertions hold for its positive
subcategory $\mathscr C_+$.

\begin{lemma}\label{lem:vol2-ambient-endomorphisms}
The category $\vM$ is presentable. Its colimits are computed
componentwise on reduced coefficients, and $\vot$ preserves
them separately. Thus $-\vot C$ has a right adjoint and
\[
 \vEnd=[C,C]_{\vot}
\]
is an ambient associative endomorphism monoid, with unit
$\eta:J\to\vEnd$ and evaluation $\vEnd\vot C\to C$.
Maps of monoids to $\vEnd$ classify coherent actions on $C$.
\end{lemma}
\begin{proof}
For positive coefficient families, the $n$th convolution power
vanishes on configurations with fewer than $n$ positions.
Bar and cobar therefore use finite filtrations at each
configuration. Normalized bar terms are $a^{\star n}[n]$
for the positive augmentation ideal $a$, and normalized cobar
terms are $c^{\star n}[-n]$ for the positive reduced coalgebra
$c$. Simplicial realization and cosimplicial totalization
give their coherent structures. Collision pullbacks preserve
these finite normalized constructions.

In $\Omega BA$, flatten a term to its $n$ original
$a$-factors and filter increasingly by $n$. Each of the
$n-1$ gaps is either inside a bar block or between cobar
factors. On the associated graded, multiplication of
original inputs disappears and the differential changes
one inside gap to an outside gap. After the suspension
signs, this is the tensor product of $n-1$ two-term identity
complexes, tensored with the coefficient word. Opposite
signs around each square give the usual tensor differential.
For $n>1$ a contraction in its first gap contracts the
product. For $n=1$ the surviving coefficient is $a$;
the empty term is $K$. The comparison is the identity
on these surviving terms.

For $B\Omega C$, use the decreasing filtration by original
$c$-factors. Coproduct splitting increases that filtration.
Its associated graded has the same gap complexes with
the direction of the cuts reversed, again leaving only
$K$ and $c$. Both filtrations end at $|I|$ on configuration
$I$, so their comparisons are equivalences. Higher coherent
products lower the first filtration and higher coproducts
raise the second; the same associated graded proof applies.
The bar and cobar adjunction retains the triangle coherences.
Thus connected coalgebras and positive augmented algebras
are equivalent as enhanced categories.

Such algebras are nonunital associative algebras in
$\mathscr C_+$. The tensor-word monad is accessible and
preserves sifted colimits, giving a presentable algebra
category \cite{LurieHA}*{Corollary~3.2.3.5}. This proves
presentability of $\vM$.
On the coalgebra side, the reduced objects are nonunital
algebras in $\mathscr C_+^{\mathrm{op}}$. Limits of algebras
are computed on their underlying objects: the operations
and their coherences are obtained by composing with each
projection. Reversing arrows proves the assertion about
coalgebra colimits.

The reduced part of a diagonal tensor is the diagonal tensor
of the reduced parts. Its pullback and external product
preserve colimits. The adjoint functor theorem
\cite{LurieHTT}*{Corollary~5.5.2.9} gives the asserted
right adjoint. Composing two evaluations and
adjointing defines multiplication on $\vEnd$. Associativity
and units are those of composition, including their higher
coherences. Adjunction transports precisely the monoid-map
identities to the coherent action identities.
\end{proof}

This lemma constructs an ambient object; it imposes no
factorization on $\vEnd$. Write
\[
 \mathscr V_C=(\vM_{/\vEnd})_{(J,\eta)/}.
\]
Then $\vAct$ is its full subcategory with factorizing
underlying object. An action $a$ will be written
$(F_a,u_a,\beta_a)$, where $\beta_a:F_a\to\vEnd$ includes
the path $\beta_a u_a\simeq\eta$.

The category $\vAct$ is accessible. For each finite ordered
partition, factorization is the condition that a natural
arrow between accessible coefficient functors be an
equivalence. The category of equivalences in an arrow
category is accessible. Pulling it back for the small set
of partitions, and then adding the pointing and slice,
gives the assertion. Choose one sufficiently large regular
cardinal $\lambda$ for these functors. The inclusion
$\vAct\to\mathscr V_C$ preserves $\lambda$-filtered
presentations: on a filtered index the diagonal to any
finite power is cofinal, so finite external products of
coefficients preserve the factorization comparisons.

Choose a small skeleton $\mathscr A_0$ of the
$\lambda$-compact generators, and let $\mathscr T_C$ be
the small full monoidal subcategory generated by it and
the unit action. Fullness retains all mapping spaces
between the selected objects. Define
\begin{equation}\label{eq:vol2-density-object}
 L_C=\operatorname*{colim}_{a\in\mathscr T_C}F_a
            \quad\text{in }\mathscr V_C,
 \qquad \lambda_a:F_a\to L_C.
\end{equation}
The unit action is initial, so this index has contractible
nerve. Its underlying colimit is consequently the colimit
in $\vM$, with the specified point and map
$\beta_L:L_C\to\vEnd$.

\begin{proposition}\label{prop:vol2-density-cocone}
The density object is a unital ambient monoid over $\vEnd$.
Its cocone extends naturally to all of $\vAct$, retaining
every mapping-space coherence, and is multiplicative for
the ordered action product $\circledast$.
\end{proposition}
\begin{proof}
Separate preservation of colimits defines multiplication by
\[
 m_L(\lambda_a\vot\lambda_b)=\lambda_{a\circledast b}.
\]
The iterated products of the index supply associativity;
its initial unit supplies both unit laws. The maps $\beta_a$
make $\beta_L$ multiplicative.

For any action $a$, the filtered slice $\mathscr A_0/a$
presents $F_a$. Thus the inclusion into $\mathscr V_C$
is the left Kan extension of its restriction to
$\mathscr A_0$. The same holds after restricting it to
$\mathscr T_C$, because that subcategory is full and
has the same slices at its objects. The universal property
of left Kan extension therefore extends the cocone,
including its higher naturality, to all actions. Passing
to compact filtered presentations in each variable extends
the multiplication identity and its coherences as well.
\end{proof}

\subsection{The section and its kernel}

For an ordered partition $P=(I_1|\cdots|I_r)$, set
\[
 X_P(F)=j_P^!F_I,\qquad
 Y_P(F)=\left(\boxtimes_{j=1}^rF_{I_j}\right)|_{U_P}.
\]
These objects lie in the stable differential-module category
on $U_P$. Internal block collisions remain in that open.
The coproduct is $\delta_{P,F}:X_P(F)\to Y_P(F)$.
For factorizing $F$ it is an equivalence; its inverse is
taken naturally in the functor category.

\begin{proposition}\label{prop:vol2-density-section}
The map $\delta_{P,L_C}$ has a section $s_P$. If
\begin{equation}\label{eq:vol2-density-kernel}
 K_{C,P}=\operatorname{fib}(\delta_{P,L_C}),\qquad
 p_P=1-s_P\delta_{P,L_C},
\end{equation}
then
\[
 X_P(L_C)\simeq Y_P(L_C)\oplus K_{C,P},\qquad
 \operatorname{cofib}(\delta_{P,L_C})\simeq K_{C,P}[1].
\]
For $h=a^{\circledast r}$ and its unit insertions $v_j:a\to h$,
put
\begin{equation}\label{eq:vol2-density-generator}
 r_{a,P}=X_P(v_1)-\delta_{P,F_h}^{-1}
             \left(\boxtimes_j(v_j)_{I_j}\right)\delta_{P,F_a}.
\end{equation}
Then $p_PX_P(\lambda_a)=X_P(\lambda_h)r_{a,P}$,
with its natural homotopy.
\end{proposition}
\begin{proof}
For independently indexed actions $a_1,\ldots,a_r$, let
$h=a_1\circledast\cdots\circledast a_r$ and let
$v_j:a_j\to h$ insert the other units. On this tuple define
\begin{equation}\label{eq:vol2-section-formula}
 s_P\big|_{(a_j)}=X_P(\lambda_h)\delta_{P,F_h}^{-1}
                   \left(\boxtimes_j(v_j)_{I_j}\right).
\end{equation}
This is a natural cocone on $\mathscr T_C^r$. The external
product preserves its separate colimits, giving a map
$Y_P(L_C)\to X_P(L_C)$. Naturality of $\delta$ makes its
composite with $\delta_{P,L_C}$ equal to the target cocone,
since $\lambda_hv_j\simeq\lambda_{a_j}$. Hence $\delta s\simeq1$.
The split fibre sequence gives the two decompositions.
Using the tuple $(a,\ldots,a)$ in
\eqref{eq:vol2-section-formula} gives
\eqref{eq:vol2-density-generator}. Compact presentations
extend the formula to every action. Internal collisions
and refinements use the same iterated coproduct coherences.
\end{proof}

\begin{theorem}\label{thm:vol2-density-criterion}
The following are equivalent: the pointed-action functor
$\Phi_C$ is representable; $L_C$ factorizes; and all
$K_{C,P}$ vanish. It suffices to test every ordered binary
block partition, with its full block coefficients.
When these conditions hold, $L_C$ is also the final
factorizing monoid action.
\end{theorem}
\begin{proof}
Binary factorization iterated inside blocks gives all
partitions. The split fibre sequence identifies
factorization with kernel vanishing.

If $L_C$ factorizes, it belongs to $\vAct$. Naturality of
the extended cocone gives
$\lambda_{L_C}\lambda_a\simeq\lambda_a$ on the small
index, hence $\lambda_{L_C}\simeq1$ by its colimit
property. Naturality on the whole space
$\vMap_{\vAct}(a,L_C)$ now homotopes its identity map
to the constant map with value $\lambda_a$. This space
is contractible, so $L_C$ is final and represents $\Phi_C$.

Conversely a final pointed action $Z$ gives a cocone
and thus a map $\rho:L_C\to Z$. The extended cocone
gives $\lambda_Z:Z\to L_C$. Naturality on the full
small index gives $\lambda_Z\rho\simeq1$. The reverse
composite is the identity by finality of $Z$. Thus
$L_C\simeq Z$ and factorizes. Finally all operation
spaces from finite products to a final object are
contractible. They give its unique coherent monoid
structure and a contractible space of monoid-action
maps into it from every factorizing monoid actor.
\end{proof}

This proof uses finality among all pointed actions.
Finality among multiplicative actions alone does not
supply the cocone from arbitrary pointed actors.

\subsection{Specified monoid extensions}

A monoid extension is $(a,M,f)$ with $M\in\vMon$ and
a pointed-action morphism $f:a\to\mathsf U M$, where
$\mathsf U$ forgets multiplication. Its morphisms include
an action map on $a$, a monoid-action map on $M$, and
a specified homotopy making their square commute.
Denote this category, with its higher coherences, by
$\mathscr E_C=\vAct^{\Delta^1}\times_{\vAct}\vMon$,
using target evaluation on the arrow category and $\mathsf U$
on the monoid category.

\begin{theorem}\label{thm:vol2-extension-cancellation}
On $\mathscr E_C$ there is a natural nullhomotopy
$p_PX_P(\lambda_a)\simeq0$. It retains the vacuum paths,
ordered coproducts and all internal collision comparisons.
\end{theorem}
\begin{proof}
For $h=a^{\circledast r}$ set
$q=\mu_M^{(r)}f^{\vot r}:h\to\mathsf U M$.
The monoid-action identity makes this a pointed-action
map in the prescribed order. Its unit paths give
$qv_j\simeq f$ for every $j$. Naturality of factorization
then gives
\[
 X_P(q)\delta_{P,F_h}^{-1}
       \left(\boxtimes_j(v_j)_{I_j}\right)\delta_{P,F_a}
 \simeq\delta_{P,M}^{-1}Y_P(f)\delta_{P,F_a}
 \simeq X_P(f).
\]
The other term of $X_P(q)r_{a,P}$ is also $X_P(f)$.
Thus $X_P(q)r_{a,P}\simeq0$. The density cocone has
$\lambda_Mq\simeq\lambda_h$, proving the assertion.

These are equations of natural transformations on
$\mathscr E_C$: the unit paths come from its coherent
monoid structures, and the factorization inverses are
natural inverses there. Their composite is therefore
a homotopy on every simplex, not a collection of
unrelated objectwise homotopies. Restriction to either
ordered collision branch and to internal block collisions
preserves the calculation. The projector's domain remains
$U_P$; no extension through cross-block collisions is asserted.
\end{proof}

The full subcategory of pointed actions which happen to
admit multiplicative structures also has maps that do
not preserve those structures. It is not the extension
category just defined. This distinction matters when
passing the nullhomotopy to a colimit.

For a small diagram $\mathscr E\to\mathscr E_C$, define
\[
 A_P^{\mathscr E}=\operatorname*{colim}_{e\in\mathscr E}
                       X_P(F_{a_e}),\qquad
 \chi_P:A_P^{\mathscr E}\longrightarrow X_P(L_C).
\]
The map is the extended density cocone. Theorem
\ref{thm:vol2-extension-cancellation} gives a nullhomotopy
of $p_P\chi_P$, so
\begin{equation}\label{eq:vol2-relative-kernel}
 \operatorname{cofib}(\chi_P)\simeq
 \operatorname{cofib}(\delta_{P,L_C}\chi_P)\oplus K_{C,P}.
\end{equation}
Indeed, put $A=A_P^{\mathscr E}$, $X=X_P(L_C)$,
$Y=Y_P(L_C)$ and $K=K_{C,P}$. In the splitting $X=Y\oplus K$, write
$\chi=(\chi_Y,\chi_K)$ and represent the specified
homotopy by $\chi_K=d_Kh+hd_A$. Our cone convention is
\[
\begin{gathered}
 \operatorname{Cone}(\chi)^n=Y^n\oplus K^n\oplus A^{n+1},\\
 d(y,k,a)=(d_Yy+\chi_Ya,d_Kk+\chi_Ka,-d_Aa).
\end{gathered}
\]
The cochain isomorphism is the shear
\begin{equation}\label{eq:vol2-cone-shear}
 (y,k,a)\longmapsto((y,a),k+h(a)).
\end{equation}
Its last differential is
$d_Kk+\chi_Ka-hd_Aa=d_K(k+ha)$, and its inverse subtracts
$ha$. In particular the retraction to the kernel includes
the chosen homotopy.

For example take $A=(k\xrightarrow{5}k)$,
$Y=(k\xrightarrow{2}k)$ and
$K'=(k\xrightarrow{(3,0)}k^2)$ in degrees $-1,0$.
Let $\chi^{-1}=(10,35)$ and $\chi^0=(4,21,0)$.
Its kernel component has homotopy $h^0=7$.
Projection to $k$ alone is not a chain map on the cone,
because of the term $21a$; the corrected projection is
$k+7a$. The nonzero class in $H^0(K')$ remains a retract.

\begin{corollary}\label{cor:vol2-extension-density}
Pointed representability holds if and only if a small
diagram of specified monoid extensions makes $\chi_P$
an equivalence for every ordered binary partition.
Equivalently, one may use a different such diagram for
each partition.
\end{corollary}
\begin{proof}
If the maps are equivalences, their cofibres vanish.
Equation~\eqref{eq:vol2-relative-kernel} makes every
$K_{C,P}$ a retract of zero, and
Theorem~\ref{thm:vol2-density-criterion} applies.
Conversely a final pointed action $Z$ has its coherent
monoid structure. Its final cocone supplies a functor
$\mathscr T_C\to\mathscr E_C$ with constant target $Z$.
The resulting $\chi_P$ is the original density colimit
on every block-disjoint open, hence an equivalence.
\end{proof}

Nonzero relative cofibres alone do not obstruct this
criterion: they may come from the first summand in
\eqref{eq:vol2-relative-kernel}. Nor is a defect before
the density colimit a class of its kernel. Take an exterior
actor with a multiplicative action, such as $(C_E,\alpha_D)$.
At two separated positions let $c_{L,i},c_{R,i}$ be the
degree-one generators in the exterior state factors of
two copies, ordered first by copy and then by position.
They give the nonzero expression
\[
 c_{L,1}c_{L,2}-c_{L,1}c_{R,2}.
\]
Multiplication sends both monomials to $c_1c_2$ and
annihilates their difference. When this multiplication
is a pointed-action map, the full density cocone
annihilates the same difference. Its action paths are
part of this conclusion.

\section{A pair surviving every pointed factorizing action}
\label{sec:vol2-pointed-sphere}

A common recipient of objects need not equalize their maps.
For the polynomial Heisenberg boundary, two prescribed maps remain
distinct after every receiving pointed factorizing action.
The vacuum imposes the restriction responsible for this obstruction:
a degree-zero operation coming from one root of a factorizing actor
must commute with the boundary-module action.

Retain $C=C_{V_\kappa}$ with $\kappa\ne0$, and put
$\mathfrak l=L_k$, $\mathfrak g=L_{V_\kappa}$ and
$A=\vU\mathfrak g$.
Write $\ell=\mathfrak l_{\{1\}}=\omega_X$ and
$G=\mathfrak g_{\{1\}}$ for the right differential modules on $X$.
Thus $J_{\{1\}}=\ell[1]$ and $C_{\{1\}}=G[1]$.
All state spaces and tensors at a fixed configuration are algebraic.

For a coalgebra $Z$, let $\overline Z=\operatorname{cofib}(K\to Z)$.
With $t=[-1]$ of degree one, the normalized cobar is
\[
 \Omega Z=K\oplus\bigoplus_{r\ge1}(t\overline Z)^{\star r},
 \qquad
 d(tc)=-t(d_Zc)+\sum_{\overline\Delta c=\sum c'\otimes c''}
                  (-1)^{|c'|+1}(tc')(tc'').
\]
Its product is concatenation. The sums are finite on every
configuration. These are the cobar and suspension conventions used
in the bar--cobar comparison above.

\subsection{An antisymmetric actor map}

Put $F_1=C_{k}$ and $F_2=C_{k^2}$ in the notation
$C_W=C_{\Lambda(W[-1])}$ of
Section~\ref{sec:vol2-nilpotent-recipients}.
Denote their exterior generators by $\varepsilon$ and
$\varepsilon_1,\varepsilon_2$, respectively.
The subscript $k$ here denotes the one-dimensional actor space;
the vacuum coalgebra remains $J=C_0$.

\begin{lemma}\label{lem:vol2-antisymmetric-actor}
There is a natural actor map $f:F_2\to F_1$, under and over $J$,
whose marked matrix and coordinate contraction are
\[
 A_f(t_1,t_2)=\tfrac12(t_2,-t_1),\qquad \Phi_f=0.
\]
Its pure linear components vanish. In a cobar representative its
pure mixed binary component is
\begin{equation}\label{eq:vol2-antisymmetric-lift}
 \tfrac12(a_mM_n-a_nM_m),\qquad
 a_m=t(s(\varepsilon m)),\quad M_n=t(sn),\quad m,n\in\mathfrak l.
\end{equation}
Every positive-weight natural action pulled back along $f$ is
naturally homotopic to the trivial action.
\end{lemma}

\begin{proof}
The matrix classification in
Theorem~\ref{thm:vol2-natural-actor-maps} already constructs an
actual map with this matrix in every weight.
An explicit composite fixes its binary representative.
Take two auxiliary odd generators $z_1,z_2$, and let $D_z$ be the
even derivation with $D_zz_1=z_2$ and $D_zz_2=0$.
The tree action extends to these graded coefficients by inserting
the Koszul sign in each permutation of the actor coefficients.

Here is a check of that extension. A nonzero term in the twisting
equation has exactly two ordinary inputs: a source bracket joins
them or joins one to an actor input, and a product uses one in each
factor. For each homogeneous coefficient input of nonnegative
degree $d$, multiply it by a product of $d$ distinct auxiliary
odd variables of degree $-1$. Use disjoint variables for different
inputs. The modified inputs have degree zero.
Extend $D_z$ trivially over these variables and apply the
fixed-point proof of Theorem~\ref{thm:vol2-profiles}.
Extracting their ordered product gives the graded twisting
identity. Permuting actor inputs gives their Koszul sign; moving
the ordinary suspension uses its degree one lower.
Thus the graded action is an actual coefficient map, with the
same finite-configuration bound.

On one $z$ input and $w$ pure auxiliary actor inputs its value is
$(D_z^wz)n\,\operatorname{pbw}(m_1\cdots m_w)$.
Act on $z_1,z_2$ by the auxiliary generator, then project
$z_1\mapsto0$, $z_2\mapsto\varepsilon$.
The resulting map $b:C_{k^3}\to F_1$ has marked matrix
\[
 B(y_1,y_2,y_3)=(y_2,0,1).
\]
Choose the inner map $a:F_2\to C_{k^3}$ with matrix
\[
 A_a(t)=
 \begin{pmatrix}
  1&0\\0&1\\-t_2/2&-t_1/2
 \end{pmatrix}.
\]
Its coordinate contraction is
$(t_1,t_2,-t_1t_2)$.
The composition law \eqref{eq:vol2-marked-composition}
gives $A_{ba}=(t_2,-t_1)/2$ and $\Phi_{ba}=0$.

For the binary calculation, relabel the intermediate generators
as $z_1,\eta,z_2$, with $\eta$ the auxiliary actor.
Write $a_z(m)=t(s(zm))$.
A lift of the mixed inner component is
\[
 L_{12}(m,n)=t\bigl(s(z_1m)s(\eta n)\bigr)
 -\tfrac12\bigl(a_{z_2}(m)M_n+a_{z_2}(n)M_m\bigr).
\]
The source input $s(\varepsilon_1\varepsilon_2p)$ lifts to
$-t(s(\eta z_1p))$. The cobar differential gives
\[
 dL_{12}
 =-t\bigl(s(\eta z_1[m,n])\bigr)
       -a_{z_1}(m)a_\eta(n)-a_\eta(n)a_{z_1}(m).
\]
This is the image of the source differential.
The primitive product correction is closed.
The relative bar contraction extends these compatible components
and keeps their binary values.
Applying $b$ sends the first term of $L_{12}$ to
$(\varepsilon m)n$ and its correction to
$-\bigl((\varepsilon m)n+(\varepsilon n)m\bigr)/2$.
This proves \eqref{eq:vol2-antisymmetric-lift}.
Both pure diagonal pairs and both pure linear components vanish.

Finally, \eqref{eq:vol2-class-substitution} identifies pullback of a
natural action class $\psi$ with $\psi(\Phi_f)$.
Its constant term is zero, so this is zero in every weight.
The Maurer--Cartan equivalence of
Theorem~\ref{thm:vol2-natural-cohomology} supplies the homotopy.
Evaluation and the relative bar comparison preserve its vacuum
path and all collision coefficients.
\end{proof}

Let $f_0:F_2\to J\to F_1$ be the augmented map.
For $D=\partial_{x_1}$ the action $\alpha_D$ has natural class $pD$.
Choose the all-weight homotopy from the preceding lemma, with the
constant homotopy for $f_0$. They give two actual morphisms
\begin{equation}\label{eq:vol2-pointed-sphere}
 (f,p_f),(f_0,p_0):
 (F_2,\mathrm{triv})\rightrightarrows(F_1,\alpha_D)
 \quad\text{in }\vAct.
\end{equation}
Equality of their first action classes alone would not have supplied
these paths.

\subsection{The operation at one root}

For a right differential module $N$ on $X$, let
$\iota_{\mathrm{diag}}N$ be its positive coefficient family, with
the direct image of $N$ on every small-diagonal branch.
For any positive family $H$ there is an equivalence
\begin{equation}\label{eq:vol2-main-diagonal-adjunction}
 \vRH_{\mathscr C}(\iota_{\mathrm{diag}}N,H)
       \simeq\vRH_X(N,H_{\{1\}}).
\end{equation}
Indeed the union of small-diagonal branches is closed on each
$\mathscr R_I$. Local cohomology there is the sum
$\bigoplus_\sigma\delta_{\sigma,*}\delta_\sigma^!H_I$.
Collision pullback preserves these unions and their local
cohomology. The resulting coefficient family is
$\iota_{\mathrm{diag}}H_{\{1\}}$, by the collision identifications.
Its counit induces an equivalence on maps from a family supported
on the branches. Restriction to one position is inverse to
$\iota_{\mathrm{diag}}$ on such supported families, proving the
formula.

Fix a pointed factorizing action $(Z,u,\gamma)$ and set
\[
 B=\Omega Z,\qquad
 \mathcal T_C=(-)\vot C,\qquad
 \Phi:\mathcal T_C(B)\longrightarrow A.
\]
Here $\Phi$ is the cobar of $\gamma$, followed by the
Chevalley--enveloping comparison.
Lemma~\ref{lem:vol2-strong-boundary-tensor} identifies
$\mathcal T_C(\Omega Z)$ with $\Omega(Z\vot C)$.
The pointing gives $\iota:\mathfrak l\to B$, and its action path
identifies $\Phi\mathcal T_C(\iota)$ with $\mathfrak g\to A$.

For an affine open $Y\subset X$ and a class
$a\in H^0\vRH_Y(\ell,B_{\{1\}})$, use
\eqref{eq:vol2-main-diagonal-adjunction} to regard $a$ as a
one-root coefficient map. Define
\[
 \Psi(a)=\Phi\mathcal T_C(a).
\]
Since $A_{\{1\}}=G$, the same adjunction identifies $\Psi(a)$
with an endomorphism of the full family $\mathfrak g$ over $Y$.
No cohomological bound on $B$ is imposed.

\begin{lemma}\label{lem:vol2-cobar-product-exchange}
On the open $U$ of two distinct roots, both ordered products
\[
 \mu_{12},\mu_{21}:
 (B_{\{1\}}\boxtimes B_{\{2\}})|_U\longrightarrow B_{\{1,2\}}|_U
\]
are equivalences, with the source order understood in each case.
For the external symmetry $\sigma$, put
$R_B=\mu_{12}^{-1}\mu_{21}\sigma$.
Every algebra map $F:B\to B'$ between such algebras satisfies
\[
 (F_1\boxtimes F_1)R_B=R_{B'}(F_1\boxtimes F_1).
\]
For $A=\vU\mathfrak g$, the exchange is the identity.
\end{lemma}

\begin{proof}
On two roots the reduced cobar has one generator
$t\overline Z_2$ and two product summands
$t\overline Z_1\boxtimes t\overline Z_1$.
Its two differential components are the coproducts of $Z$,
with the displayed suspension signs.
Both are equivalences on $U$.
Quotienting by either product summand leaves the cone of the other
coproduct equivalence. Thus both product inclusions are equivalences,
even for unbounded complexes.
The exchange identity follows after composition with $\mu'_{12}$
from preservation of the two products.
In $A$, the difference of the two generator products is the
chiral bracket, supported on the diagonal. Hence its open exchange
is the identity.
\end{proof}

Write $\mu_{\mathfrak g}:\mathfrak g\star\mathfrak g\to\mathfrak g$
for the coefficient bracket.
An endomorphism $M$ of the module $\mathfrak g$ obeys
$M\mu_{\mathfrak g}=\mu_{\mathfrak g}(\vid\star M)$.

\begin{lemma}\label{lem:vol2-vacuum-module-map}
Every one-root operation $M=\Psi(a)$ satisfies
\begin{equation}\label{eq:vol2-vacuum-module-map}
 M\mu_{\mathfrak g}
 =\mu_{\mathfrak g}(\vid\star M)
 =\mu_{\mathfrak g}(M\star\vid)
\end{equation}
on both ordered branches. If it is also a coefficient derivation,
then $M=0$.
\end{lemma}

\begin{proof}
All maps are derived coefficient maps over $Y$.
On the two-root scheme write $j:U\hookrightarrow\mathscr R_2$
for the disjoint open and $i$ for the union of the two closed
diagonal branches. Put $R=\ell\boxtimes\ell$, pulled back from
$Y^2$. Both convolution summands have source $j_*j^*R$.
The difference $\kappa_a=[\iota,a]$ of their ordered products obeys
\[
 \Phi\mathcal T_C(\kappa_a)
       =\mu_{\mathfrak g}(\vid\star M)
\]
by the envelope relation and the chosen unit path.

The commutator may have an open part. Define on $U$
\[
 w=(R_B-\vid)(\iota\boxtimes a),\qquad
 P(w)=\mu_{12}j_*(w):j_*j^*R\longrightarrow B_2.
\]
The second formula uses the global product from the first
convolution summand. Thus $\kappa_a+P(w)$ restricts to zero on $U$.
The exchange identity for $\Phi$, and the strong symmetric
monoidal comparison for $\mathcal T_C$, give
\[
 \Phi\mathcal T_C(P(w))=0.
\]
Indeed the image of $w$ is killed by $R_A-\vid$.
The open projection formula carries this zero map and its
nullhomotopy through $j_*$ and multiplication.
No primitive is extended from the open.

Put $Q=\operatorname{cofib}(j_!j^*R\to j_*j^*R)$.
Open adjunction lifts $\kappa_a+P(w)$ through $Q$.
The source $R$ is regular holonomic, so $j_!$ is defined for this
source regardless of the unbounded target.
The two localization triangles give
\[
 i_*i^*R\longrightarrow Q\longrightarrow i_*i^!R[1].
\]
On each branch $i^!R=\ell[-1]$ and $i^*R=\ell[1]$.
The exact sequence of principal parts
$\omega\to\omega(*0)\to\delta_0$ fixes the first shift;
duality fixes the second.
These are the right-module conventions of
\cite{BD04}*{\S\S2.1.1--2.1.3}.
The map $j_*j^*R\to Q\to i_*\ell$ is the vacuum
principal-part bracket.

The closed projection formula and $i^!C_2=C_1$ give
\[
 \mathcal T_C(i_*i^*R)=i_*G[1],\qquad
 \mathcal T_C(i_*i^!R[1])=i_*G.
\]
The finite PBW filtration puts $A_2$ in the right-module heart.
Its quotients are open and closed direct images of degree-zero
external products; these images are exact on the affine collision
charts. Also $i^!A_2=G$.
Consequently
\[
 \operatorname{Hom}(i_*G[1],A_2)
   =\operatorname{Ext}^{-1}(G,G)=0.
\]
After applying $\mathcal T_C$ and $\Phi$, the lifted map from $Q$
therefore factors on each branch through an endomorphism $N$ of $G$.
It follows that
\[
 \mu_{\mathfrak g}(\vid\star M)=N\mu_{\mathfrak g}.
\]

Insert the vacuum in the first slot. The bracket becomes
$\mathrm{can}:j_*j^*(\ell\boxtimes G)\to i_*G$,
and the last equality becomes $(N-M)\mathrm{can}=0$.
Precomposition by $\mathrm{can}$ is injective on degree-zero
maps to $A_2$: its kernel is an image of
$\operatorname{Hom}((\ell\boxtimes G)[1],A_2)=0$,
again by the right-module heart.
Thus $N=M$. Skew symmetry gives the identity with $M$ in the
first slot, proving \eqref{eq:vol2-vacuum-module-map}.

If $M$ is a derivation, subtract its second-slot module identity
from the derivation identity.
The result is $\mu_{\mathfrak g}(M\star\vid)=0$.
Inserting the vacuum in the second slot and using the same
principal-part injectivity gives $M=0$.
\end{proof}

\begin{corollary}\label{cor:vol2-one-root-scalars}
Every $\Psi(a)$ over a nonempty affine open is a constant scalar
multiple of the identity.
\end{corollary}

\begin{proof}
The module identity implies commutation with every state mode.
Normal transfer can be written with the normal derivative in the
first coordinate; the matrix of $M$ on the second coordinate then
acquires no derivatives.
Put $c=M\mathbf1$. The regular coefficient of $J(z)$ multiplied by $(z-w)^{-n}$
detects the creation mode $J_{-n}$ in the principal-part module
identity. Commutation with these modes gives
$M(f)=fc$ for every polynomial state.
Commutation with $J_n=\kappa n\partial_{x_n}$ for $n>0$ gives
$\partial_{x_n}c=0$. Thus $c$ is a scalar function of the root.
Horizontality makes it constant.
\end{proof}

Both kinds of modes are necessary. The derivation
$D=\partial_{x_1}$ commutes with all annihilation modes, but
\[
 (DJ_{-1}-J_{-1}D)\mathbf1=\mathbf1.
\]
It is not a module map.

\subsection{A singleton class that no recipient kills}

Normalize the fibre of a right differential module at $q\in X(k)$ by
\[
 N_q=i_q^!N[1].
\]
Then $\ell_q=k$ and $G_q=V_\kappa$ lie in degree zero.
This is a derived restriction, so its induced maps on coefficient
complexes preserve boundaries.

\begin{lemma}\label{lem:vol2-affine-kunneth}
For a nonempty affine open $Y\subset X$, a complex $N$ of
differential modules on $Y$ and a complex $W$ of $k$-vector spaces,
\[
 \vRH_Y(\ell,N\otimes_kW)
 \simeq\vRH_Y(\ell,N)\otimes_kW.
\]
No bound on either complex is required. Every cohomology class
has finite support in the algebraic K\"unneth direct sum.
\end{lemma}

\begin{proof}
The unshifted left--right conversion takes $\ell$ to $\mathcal O_Y$.
For the affine coordinate $z$, the Spencer resolution is
\[
 0\longrightarrow D_Y
 \xrightarrow{P\mapsto P\partial_z}D_Y
 \xrightarrow{P\mapsto P(1)}\mathcal O_Y\longrightarrow0.
\]
Affineness makes it a two-term resolution by finite free modules
for the derived Hom in question. Its Hom complex commutes with
tensoring by any complex $W$. Over the field $k$, the cohomology
of an algebraic tensor product is the direct sum of the tensor
products of cohomology groups. Every element has finite support.
\end{proof}

Let $D_0\in\mathfrak d_V$ be a coefficient derivation, and let
$\alpha$ be any action of $F_1$ with first-weight class
$[\gamma_{D_0}^{(1)}]$ in
Theorem~\ref{thm:vol2-first-weight-classification}.
This includes the natural action $\alpha_{D_0}$.
For a receiving morphism
\[
 h:(F_1,\alpha)\longrightarrow(Z,\gamma),
\]
let $h_0$ be its twisting component on the pure singleton
$s(\varepsilon m)$, $m\in\mathfrak l$.
It defines a class
\[
 [h_0]\in H^1\vRH_X(\ell,B_{\{1\}}).
\]
This class forgets equivariance by $\mathfrak l$.
Nonvanishing of an equivariant extension would not by itself
prove its nonvanishing here.

\begin{theorem}\label{thm:vol2-plain-singleton}
If $D_0\ne0$, then $[h_0]|_Y\ne0$ on every nonempty affine
open $Y\subset X$, for every receiving pointed factorizing action.
\end{theorem}

\begin{proof}
Suppose $[h_0]|_Y=0$. Resolve the pointed arrow before taking
its relative convolution, retaining the full coefficient
differential. A primitive of $h_0$ defines a positive-actor-weight
gauge. Its weight-one equation replaces $h_0$ by zero exactly;
higher weights change by the usual gauge series, finite at each
configuration. This gives $h'\simeq h$ under $J$ with $h'_0=0$.
Composition transports the specified action path.

The first mixed component of $h'$ has one ordinary input and
one actor input. It is a closed degree-zero cochain $c_1$.
It need not be supported on collisions.
The normalized bar formula for $\mathcal T_C(h')$ on these inputs
has the term $\mathcal T_C(c_1)$ and products involving
$\mathcal T_C(h'_0)$. Those products vanish. Its composite
with the action therefore has mixed component
$\Phi\mathcal T_C(c_1)$.

Choose $q\in Y(k)$ and put $Y'=Y\setminus\{q\}$.
Such a choice exists because $k$ is infinite and $Y$ is a nonempty
open of $\mathbb A^1_k$.
Keep the actor root at $q$ and let the ordinary root vary over $Y'$.
The first ordered product and normalized fibre identify the target
on $Y'\times\{q\}$ with $B_{\{1\}}|_{Y'}\otimes_kB_q$.
Lemma~\ref{lem:vol2-affine-kunneth} writes the class of $c_1$
as a finite sum with bidegrees $(r,-r)$.
After the action, its second factors lie in
\[
 H^{-r}\vRH_k(V_\kappa,V_\kappa)=0\qquad(r\ne0).
\]
Thus the resulting operation has the form
\[
 \sum_{i=1}^m M_i\otimes N_i,\qquad
 M_i=\Psi(a_i),\quad N_i=\Psi_q(b_i),
\]
where $a_i$ and $b_i$ have degree zero and $\Psi_q$ is the
normalized point action. The module maps $M_i$ act on the moving
ordinary-root factor.

The action path identifies the first-weight class of this composite
with $[\gamma_{D_0}^{(1)}]$.
The two-output part of the proof of
Theorem~\ref{thm:vol2-first-weight-classification} identifies such
classes with actual module cocycles, without boundaries.
Indeed open adjunction kills the zero Chevalley column in every
degree, and total degree one leaves only its first column in
coefficient degree zero. This argument also holds on $Y$.
Restricting the resulting equality of maps gives
\begin{equation}\label{eq:vol2-singleton-finite-sum}
 D_0|_{Y'}\otimes\vid_{V_\kappa}
       =\sum_{i=1}^m M_i\otimes N_i.
\end{equation}

Algebraic tensors of endomorphisms inject into external
endomorphisms. For the finite list of second factors, finitely
many vector evaluations and linear functionals separate every
linearly independent sublist; applying them recovers each first
factor. Choose a functional $\lambda$ on
$\operatorname{span}_k\{\vid,N_1,\ldots,N_m\}$ with
$\lambda(\vid)=1$. Equation~\eqref{eq:vol2-singleton-finite-sum}
then gives
\[
 D_0|_{Y'}=\sum_i\lambda(N_i)M_i.
\]
Lemma~\ref{lem:vol2-vacuum-module-map} makes the right side a
module map. Since $D_0$ is a derivation, that lemma forces
$D_0|_{Y'}=0$.
The coefficient module $\mathcal O_X\otimes_kV_\kappa$ is
torsion-free, and every state has an algebraic image.
Localization is faithful on each such image, so $D_0=0$ on $X$,
a contradiction.
\end{proof}

\subsection{Homotopy invariance on two roots}

The obstruction must survive homotopies of coalgebra maps,
including those with nonzero intermediate singleton components.
Use compatible differential graded models of the full coefficient
mapping complexes, with the pointed arrow resolved, and write
$\partial$ for their differential.
The two pure inputs $s(\varepsilon_1m)$ and
$s(\varepsilon_2n)$ have degree zero. Their bidegrees record the
numbers of $\varepsilon_1$ and $\varepsilon_2$.
Let $j$ denote restriction to the two-root disjoint open and put
$P_{1,1}=\mathfrak l\star\mathfrak l$, keeping the labels.

\begin{lemma}\label{lem:vol2-binary-homotopy}
Suppose maps $\phi_0,\phi_1:F_2\to Z$ under $J$ have twisting
representatives with zero pure singleton actor components.
If they are homotopic under $J$, their restricted pure binary
components give the same class in
\[
 H^1\vRH(j^*P_{1,1},j^*B).
\]
\end{lemma}

\begin{proof}
Represent the relative homotopy by $\tau(u)+c(u)\,du$,
where $\tau$ and $c$ have degrees one and zero.
Polynomial interval forms suffice in the positive filtration,
which is finite on each configuration.
On a pure singleton input the source differential and reduced
coproduct vanish. Hence
\[
 \tau_{1,0}=\partial E,\quad \tau_{0,1}=\partial F,
 \qquad E(u)=\int_0^u c_{1,0}(v)\,dv,
 \quad F(u)=\int_0^u c_{0,1}(v)\,dv.
\]
We have $E(0)=F(0)=0$ and
$\partial E(1)=\partial F(1)=0$.
On the disjoint open the source bracket vanishes. The mixed
interval equation, with a prime denoting $u$-derivative, is
\[
 \tau'_{1,1}=\partial c_{1,1}
             +[\partial E,F']+[\partial F,E'].
\]
The brackets use both labelled external products followed by
the target multiplication. Since $E$ and $F$ have degree zero,
the graded Leibniz rule gives
\begin{equation}\label{eq:vol2-binary-boundary}
 [\partial E,F']+[\partial F,E']
       =\partial[E,F']-\frac{d}{du}[E,\partial F].
\end{equation}
Therefore
\[
 \tau_{1,1}(1)-\tau_{1,1}(0)
 =\partial\int_0^1\bigl(c_{1,1}+[E,F']\bigr)\,du
       -[E,\partial F]\big|_0^1.
\]
The endpoint term is zero. At both endpoints the
Maurer--Cartan equation makes the restricted binary term closed,
because the singleton terms vanish. Their difference is thus a
boundary in the stated full coefficient complex.
\end{proof}

The commutator correction need not be zero as a cochain.
For $E=u(1-u)E_0$ and $F=u^2(1-u)F_0$, with $E_0,F_0$
of degree zero in a differential graded algebra,
\[
 \int_0^1[E,F']\,du=\frac1{60}[E_0,F_0].
\]
Its differential can be nonzero; the lemma uses its exactness.

\begin{theorem}\label{thm:vol2-pointed-sphere}
For every morphism
$h:(F_1,\alpha_D)\to(Z,\gamma)$ in $\vAct$, the composites
$hf$ and $hf_0$ from \eqref{eq:vol2-pointed-sphere} are not
homotopic under $J$.
Thus $\vAct$ is not filtered, and $\Phi_C$ is not representable
on pointed factorizing coalgebras for the polynomial Heisenberg
boundary at nonzero level.
\end{theorem}

\begin{proof}
Both composites have zero pure singleton actor components.
On two distinct roots the source cobar differential gives
\[
 d\,t\bigl(s(\varepsilon n)s(m)\bigr)
 =-t\bigl(s(\varepsilon n)\bigr)t\bigl(s(m)\bigr)
   +t\bigl(s(m)\bigr)t\bigl(s(\varepsilon n)\bigr).
\]
Applying $\Omega h$ reorders $h_0(n)\iota(m)$ into
$\iota(m)h_0(n)$ modulo a boundary.
The boundary comes from the cocommutative source; no centrality
of the vacuum in $Z$ is assumed.

Fix $q\in X(k)$ and put $Y=X\setminus\{q\}$.
Restrict to $Y\times\{q\}$ by the normalized derived fibre.
In the first product order, the binary class of $hf$ is
\begin{equation}\label{eq:vol2-sphere-binary-class}
 \frac12\bigl([h_0]|_Y\otimes[\iota_q]
                  -[\iota]|_Y\otimes[h_0]_q\bigr).
\end{equation}
Lemma~\ref{lem:vol2-affine-kunneth} puts the two terms in the
distinct K\"unneth summands of bidegrees $(1,0)$ and $(0,1)$.
The first is nonzero by Theorem~\ref{thm:vol2-plain-singleton}
and $\Psi_q(\iota_q)=\vid_{V_\kappa}$.
The second cannot cancel it and need not itself be nonzero.
The composite $hf_0$ has zero binary class.
Lemma~\ref{lem:vol2-binary-homotopy} excludes a homotopy under
$J$, and hence a homotopy of pointed actions with its action
two-cell.

A filtered infinity-category equalizes every parallel pair
after some receiving morphism. The pair
\eqref{eq:vol2-pointed-sphere} violates that condition.
A representing object for $\Phi_C$ would be a final object of
$\vAct$; its mapping space from $(F_2,\mathrm{triv})$ would
be contractible. That too would identify the two composites,
which is impossible.
\end{proof}

\begin{corollary}\label{cor:vol2-independent-unit-products}
For every such $h$, the two classes represented on the disjoint
open by $h_0(m)\iota(n)$ and $h_0(n)\iota(m)$ are linearly
independent in $H^1\vRH(j^*P_{1,1},j^*B)$.
\end{corollary}

\begin{proof}
Write the classes as $c_1,c_2$, with the root labels fixed.
Fix $q\in X(k)$ and put $Y=X\setminus\{q\}$.
Suppose $a c_1+b c_2=0$. On $Y\times\{q\}$ the
$(1,0)$ K\"unneth component is
$a[h_0]|_Y\otimes[\iota_q]$, which is nonzero unless $a=0$.
Now restrict to $\{q\}\times Y$ and take the moving root first
in the product order. The same component is
$b[h_0]|_Y\otimes[\iota_q]$, so $b=0$.
The source cobar boundary used in the theorem permits both
reorderings. Thus neither centrality of the recipient vacuum nor
nonvanishing of $[h_0]_q$ is required.
\end{proof}

Multiplying the marked matrix of $f$ by $2\lambda$, for
$\lambda\in k^\times$, replaces $1/2$ in
\eqref{eq:vol2-sphere-binary-class} by $\lambda$.
The coordinate contraction remains zero, so the same construction
and obstruction apply.

The theorem concerns pointed actions with all pointed maps.
For the additive exterior multiplications, $f$ is not a monoid map.
Its first marked form is $(t_2\,dt_1-t_1\,dt_2)/2$.
Pulling it back by addition and subtracting its values on the two
factors leaves
\[
 \tfrac12(s_2\,dt_1-s_1\,dt_2+t_2\,ds_1-t_1\,ds_2),
\]
which is nonzero in the first marked-matrix invariant.
Thus the displayed pair does not decide finality in $\vMon$,
where comparison maps must preserve multiplication.

\subsection{How the ambient endomorphism object kills the class}

The object $\vEnd$ of Lemma~\ref{lem:vol2-ambient-endomorphisms}
receives every action, without a factorization requirement.
Its failure to factorize can be seen directly in its cochains.
Put $c=\overline C$ and $a=\overline{\Omega C}$, and let
$\underline{\vRH}^{!}$ be the internal Hom for $\vot$ on positive
coefficient families. There is an equivalence
\begin{equation}\label{eq:vol2-internal-action-cochains}
 \Omega\vEnd\simeq K\oplus\underline{\vRH}^{!}(c,a).
\end{equation}
The differential on the right is the internal mapping
differential. Its cup product uses the coproduct of $C$,
the interchange map, evaluation, and multiplication in $\Omega C$.

To prove the underlying equivalence, test against the free augmented
algebra $T_\star(V)$ on a positive coefficient family $V$.
Its bar contracts to $K\oplus sV$, with zero reduced coproduct:
the same cut contraction as in
Lemma~\ref{lem:vol2-ambient-endomorphisms} leaves only the
terms of tensor length one. Tensoring with $C$ leaves the primitive reduced
coefficient $sV\vot c$. The bar--cobar and internal-Hom
adjunctions identify maps successively with
\[
 \vMap_{\mathscr C}(V,\overline{\Omega\vEnd})
 \simeq\vMap_{\vM}\bigl((K\oplus sV)\vot C,C\bigr)
 \simeq\vMap_{\mathscr C}\bigl(V,
                    \underline{\vRH}^{!}(c,a)\bigr).
\]
Yoneda gives the equivalence and its differential.
For the product, take strict representatives and put
$H=K\oplus\underline{\vRH}^{!}(c,a)$ with the stated cup product.
Let $\pi:\overline{\vB H}\to\underline{\vRH}^{!}(c,a)$ be the
bar twisting cochain. Evaluation gives
\[
 \tau=\operatorname{ev}(\pi\vot1):
                 \overline{\vB H}\vot c\longrightarrow a.
\]
Naturality of interchange identifies the quadratic term in the
twisting equation for $\tau$ with the evaluation of the cup-product
term for $\pi$. Evaluation is a chain map, so the linear terms
agree as well. Thus $\tau$ defines $\vB H\vot C\to C$.
Adjunction and cobar give $H\to\Omega\vEnd$.
On a map $V\to\underline{\vRH}^{!}(c,a)$, its adjoint is exactly
evaluation of that map. The underlying ideal map is therefore
the identity under the equivalence just proved. This proves the
algebra equivalence in \eqref{eq:vol2-internal-action-cochains}.

\begin{proposition}\label{prop:vol2-ambient-kills-singleton}
For every coefficient derivation $D_0\in\mathfrak d_V$, the
classifying map $\beta_{\alpha_{D_0}}:F_1\to\vEnd$ has zero
singleton actor component. For every receiving morphism
$h:(F_1,\alpha_{D_0})\to(Z,\gamma)$, the classifying map
$\beta_\gamma:Z\to\vEnd$ sends $[h_0]$ to zero.
The ambient endomorphism object $\vEnd$ is not factorizing.
\end{proposition}

\begin{proof}
The singleton classifying map is adjoint to the singleton action.
The twisting cochain for $\alpha_{D_0}$ needs one ordinary input;
a pure singleton actor input has none. Its image is therefore
zero. The specified action path gives
$\beta_\gamma h\simeq\beta_{\alpha_{D_0}}$ under $J$.
Restriction to one position proves the assertion about $[h_0]$.

Equation~\eqref{eq:vol2-internal-action-cochains} supplies a
degree-zero internal cochain $a_D:\mathfrak l\to\Omega\vEnd$
whose one-root action is $D$.
At a singleton, its adjunction is precisely that between
$\ell\vot C_1=G$ and $D:G\to G$;
\eqref{eq:vol2-main-diagonal-adjunction} extends it to the full
coefficient family. If $\vEnd$ factorized,
Lemma~\ref{lem:vol2-vacuum-module-map} would make $D$ a module
map, contradicting $(DJ_{-1}-J_{-1}D)\mathbf1=\mathbf1$.
\end{proof}

For any nonzero coefficient derivation $D_0$, the same adjunction
gives $a_{D_0}:\mathfrak l\to\Omega\vEnd$.
Its vacuum commutator $[\iota,a_{D_0}]$ is nonzero on the
two-root disjoint open.
If it vanished there, the commutator itself would lift through
$Q$ in the proof of Lemma~\ref{lem:vol2-vacuum-module-map}.
The supported part of that proof would give
$\mu_{\mathfrak g}(\vid\star D_0)=N\mu_{\mathfrak g}$.
The first vacuum slot forces $N=D_0$; the derivation identity
and the second vacuum slot then force $D_0=0$.
Once open vanishing is assumed, this argument uses no
factorization of the target.
The nonzero open commutator becomes zero after the action,
whose bracket image is supported on the diagonal.
This explains how $\vEnd$ can carry the mixed derivation
while killing the singleton class.

Finally, Theorems~\ref{thm:vol2-density-criterion} and
\ref{thm:vol2-pointed-sphere} imply that $L_C$ is not factorizing:
at least one of its two-block kernel coefficients $K_{C,A|B}$ is
nonzero. This consequence does not yet exhibit a representative
of such a coefficient in the full density bar. That bar also
retains maps capable of killing defects of individual objects.

\section{Multiplication and a surviving commutator}
\label{sec:vol2-monoid-obstruction}

The pointed sphere did not use a multiplication on its source
actor. To obstruct a final multiplicative action, we need two
monoid-action maps. A trivially acting submonoid supplies them:
its inclusion and the map through the unit have the same action.
The multiplication of a larger actor can nevertheless distinguish
these maps after every receiving action.

Work over $\mathbb C$. Let $Q$ be an algebraic unital vertex algebra
in cohomological degree zero, with translation $T$, and let
$C_Q=\vCE(L_Q)$ on the full ordered collision carrier.
A vertex derivation is a linear map $D:Q\to Q$ satisfying
\[
 D\mathbf1=0,\qquad DT=TD,\qquad
 D(x_{(n)}y)=(Dx)_{(n)}y+x_{(n)}Dy\quad(n\in\mathbb Z).
\]
We retain the category of multiplicative actions defined in
Section~\ref{sec:vol2-action-carrier}. Receiving actors may be
unbounded. They need not have an augmentation to $J$.

\subsection{An exterior monoid with a trivial direction}

Let $A_0=\mathbb C a\oplus\mathbb C b$ have multiplication
\[
 a^2=a,\qquad ab=ba=b,\qquad b^2=0.
\]
The unit $a$ belongs to this auxiliary algebra; its exterior
generator is distinct from the vacuum. Set
\[
 \mathsf F=C_{\Lambda(A_0[-1])},\qquad
 \mathsf G=C_{\Lambda(\mathbb C a[-1])},\qquad
 \mathsf H=C_{\Lambda(\mathbb C b[-1])}.
\]
The exterior generators have degree one, zero differential and
zero translation. For $x\in A_0$ let $L_x$ be multiplication by
$x$ and put
\[
 R(x)=e^{L_x},\qquad
 M(x)=\sum_{n\geq0}\frac{L_x^n}{(n+1)!}.
\]
The constant term of $M$ is the identity. Define the marked
multiplication matrix
\begin{equation}\label{eq:vol2-dual-number-product}
 (\mathcal L(x,y)\mid\mathcal R(x,y))
 =\bigl(M(x+y)^{-1}M(x)\mid
        M(x+y)^{-1}R(x)M(y)\bigr).
\end{equation}

\begin{lemma}\label{lem:vol2-dual-number-monoid}
This matrix gives a coherent unital monoid on $\mathsf F$ under
and over $J$, with coordinate contraction $x+y$.
The two inclusions
$i:\mathsf G\to\mathsf F$ and $j:\mathsf H\to\mathsf F$
are monoid maps. The induced multiplication on $\mathsf H$ is
additive. All these maps preserve the auxiliary grading
$\deg a=0$, $\deg b=1$.
\end{lemma}
\begin{proof}
The law
\[
 (x,U)(y,V)=(x+y,U+R(x)V)
\]
is associative because $R(x+y)=R(x)R(y)$.
Changing fibre coordinates by $U=M(x)u$ gives
\eqref{eq:vol2-dual-number-product}. Its three associativity
identities are \eqref{eq:vol2-multiplication-blocks}, and its
unit identities follow from $M(0)=I$. In $A_0$,
\[
 M(x)x=e^x-a,\qquad
 e^{x+y}-a=(e^x-a)+e^x(e^y-a).
\]
Thus the diagonal $u=x$ is a homomorphism, proving the coordinate
contraction. The $a$- and $b$-lines are preserved; on the $b$-line,
$L_{sb}b=0$, so both multiplication blocks are the identity.
The algebra product preserves auxiliary degree, hence so do
$R$, $M$ and $M^{-1}$.

Theorem~\ref{thm:vol2-natural-actor-maps} realizes these matrices
and their identities, with a contractible space of coherent
choices. The multigraded contractions in its proof preserve the
auxiliary grading. At $m$ positions, exterior degree bounds the
source actor weight by $2m$. The formal series therefore evaluate
as finite sums on the original coefficient diagrams.
\end{proof}

The zero auxiliary degree of $a$ prevents the use of positive
grading alone to remove action loops. Prescribing the full
$\mathsf G$-action removes those loops instead.

\begin{lemma}\label{lem:vol2-relative-dual-number-action}
For every vertex derivation $D$, there is a coherent unital
$\mathsf F$-action $\alpha$ on $C_Q$ whose restriction to
$\mathsf G$ is trivial and whose restriction to $\mathsf H$ is
the single-derivation action $\alpha_D$.
With these comparisons, $i$ and $j$ are monoid-action maps.
\end{lemma}
\begin{proof}
First use one formal derivation. For $r$ actor copies let
$t_\nu,s_\nu$ be coordinates dual to $a,b$. Give them auxiliary
weights $0,-1$ and give $D$ weight $1$. Take natural cochains
of weight zero which vanish when all $b$-inputs are absent.
The scalar-bar and Koszul projection of
Theorem~\ref{thm:vol2-natural-cohomology} has target
\begin{equation}\label{eq:vol2-relative-monoid-model}
 \left[(s_1,\ldots,s_r)
 \mathbb C[[t_1,s_1,\ldots,t_r,s_r]]
       \otimes\Lambda(\mathbb C D[-1])\right]^{\mathbb G_m}.
\end{equation}
Completion is in total actor weight. Restriction to the
$a$-subspace has a splitting, and taking auxiliary weight zero
is exact in each finite total weight and derivative order.
The projection is consequently surjective with weightwise
acyclic kernel on this relative complex.

The target has no nonpositive degrees: a scalar monomial in
the displayed ideal has negative auxiliary weight. Its degree-one
part consists of $\sum_\nu s_\nu f_\nu(t_1,\ldots,t_r)D$.
Relative Maurer--Cartan lifting therefore makes each corresponding
component contractible. The same assertion holds with prescribed
unit faces and restrictions to the $b$-lines. Indeed the deletion
projections for actor copies and the two exterior directions
commute and split. Inclusion--exclusion extends compatible data
on their union; the residual kernel is a sum of the same acyclic
multigraded kernels. Polynomial forms relative to parameter
boundaries then supply the required fillings.

Choose the class $sD$ on one copy. Pullback by multiplication
has class $(s_1+s_2)D$, since the coordinate contraction is
addition. The composite action has the same class by
\eqref{eq:vol2-action-BCH}, or directly by
$e^{s_1D}e^{s_2D}=e^{(s_1+s_2)D}$.
Contractibility supplies its multiplication path, relative to
the prescribed line actions and units. Each subsequent module
associahedron has boundary in the component with class
$(s_1+\cdots+s_r)D$. Relative contractibility fills it and all
parameterized comparison boundaries.

Finally evaluate on $Q$. The strong symmetric monoidal functor
$\mathcal T_Q=(-)\vot C_Q$ of
Lemma~\ref{lem:vol2-strong-boundary-tensor} sends the vacuum
Lie object to $L_Q$, with its bracket. It therefore identifies
$\mathsf F\vot C_Q$ with
$\vCE(L_Q\otimes\Lambda(A_0[-1]))$ and commutes with envelopes
and bars. The infinitesimal vertex automorphism $1+uD$ over
$\mathbb C[u]/(u^2)$ shows that this comparison carries $D$ to
the coefficient derivation, retaining collisions and normal
transfers. It carries the universal polynomials and their
parameter homotopies as well. At each configuration, actor
weight is bounded and derivative order equals the $b$-count.
Every evaluated operation is finite.
\end{proof}

\subsection{A commutator detected on two roots}

Let $\mathsf S=C_{\Lambda(\varepsilon_1,\varepsilon_2)}$.
Map it to $\mathsf F^{\vot4}$ by
\[
 \varepsilon_1\longmapsto a^{(1)}-a^{(3)},\qquad
 \varepsilon_2\longmapsto b^{(2)}-b^{(4)},
\]
and multiply the four factors in that order. Call the resulting
pointed map $w:\mathsf S\to\mathsf F$. Define the formal series
$\chi(t)=\sum_{n\geq0}t^n/(n+1)!$.

\begin{lemma}\label{lem:vol2-monoid-commutator}
There is a natural pointed map $f:\mathsf S\to\mathsf H$ with
\begin{equation}\label{eq:vol2-monoid-commutator}
 w\simeq jf,\qquad A_f(t,s)=(-s\chi(t)\mid t\chi(t)).
\end{equation}
Its coordinate contraction and singleton actor components vanish.
Its first marked form is $t\,ds-s\,dt$.
\end{lemma}
\begin{proof}
In the affine coordinates used above, an $a$-line element with
coordinate $t$ and fibre $ra$ has linear part $e^t$ and
translation $\chi(t)ra$. A $b$-line element with coordinate $s$
and fibre $qb$ has linear part multiplication by $a+sb$ and
translation $qb$. Negating coordinate and fibre gives inversion,
since $M(-x)=R(-x)M(x)$. The linear parts commute. Their affine
commutator has linear part one and translation
\[
 -s\chi(t)rb+(e^t-1)qb=\chi(t)(tq-sr)b.
\]
Its coordinate is zero, where the frame is the identity.
The marked-map classification now gives \eqref{eq:vol2-monoid-commutator}
and its pointed homotopy. Finally $A_f(t,s)(t,s)^T=0$,
$A_f(0,0)=0$ and $\chi(0)=1$ prove the remaining assertions.
\end{proof}

The singleton detector extends from the Heisenberg algebra to
$Q$ without a scalar-endomorphism classification.
Let $h:(\mathsf H,\alpha_D)\to(Z,\beta)$ be any receiving
pointed factorizing action and put $B=\Omega Z$.
Write $h_0$ for its pure singleton twisting component and
$\iota:L_{\mathbb C}\to B$ for its pointing.

\begin{lemma}\label{lem:vol2-vertex-singleton}
If $D\ne0$, then $[h_0]|_Y\ne0$ in
$H^1\vRH_Y(\ell,B_{\{1\}})$ for every nonempty affine
open $Y\subset X$.
\end{lemma}
\begin{proof}
The proof of Lemma~\ref{lem:vol2-vacuum-module-map} applies
with $G=\omega_X\otimes Q$ and $\mathfrak g=L_Q$.
Here are the hypotheses it uses. Factorization makes both
ordered cobar products invertible on distinct roots. The open
commutator correction has zero action by the exchange identity
and open projection. Its remaining supported part factors
through the vacuum localization cofiber, whose terms on each
branch are $\ell[1]$ and $\ell$. Tensoring with $C_Q$ replaces
them by $G[1]$ and $G$. The PBW envelope of $L_Q$ lies in the
right-module heart on each affine collision chart, with
extraordinary diagonal restriction $G$. Negative Ext therefore
kills the first term and makes the principal-part vacuum map
injective on degree-zero maps. Inserting the vacuum in both
slots gives
\[
 M\mu_{L_Q}=\mu_{L_Q}(1\star M)=\mu_{L_Q}(M\star1)
\]
for every plain degree-zero coefficient operation $M$.
Such an operation that is also a derivation is zero.

If $[h_0]|_Y=0$, the relative gauge used in
Theorem~\ref{thm:vol2-plain-singleton} removes $h_0$.
Separate the moving ordinary root on $Y'=Y\setminus\{q\}$
from the actor root at $q\in Y$. The finite Spencer comparison
of Lemma~\ref{lem:vol2-affine-kunneth} applies to arbitrary
receiving complexes. Since $Q$ is in degree zero, only the
bidegree $(0,0)$ survives after its action. The action path gives
\[
 D|_{Y'}\otimes1_Q=\sum_{\nu=1}^m M_\nu\otimes N_\nu.
\]
This is equality of two-output cocycles. Resolve the coefficient
maps in nonnegative degrees before taking the Chevalley complex.
Open adjunction annihilates its zero column; total degree one
then has only column one in coefficient degree zero, with no
incoming boundary. Each $M_\nu$ satisfies the module
identities just proved. Finite vector evaluations separate the
span of $1_Q,N_1,\ldots,N_m$. Applying a functional taking
$1_Q$ to $1$ makes $D|_{Y'}$ a linear combination of these
module maps. It is a derivation as well, so it vanishes.
Localization is faithful on the free algebraic state module
$\mathcal O_X\otimes Q$. Hence $D=0$, a contradiction.
\end{proof}

\begin{theorem}\label{thm:vol2-monoid-nonfinality}
If $Q$ has a nonzero vertex derivation, its category of coherent
unital multiplicative actions on ordinary factorizing coalgebras
has no final object. There are parallel monoid-action maps
which no receiving monoid action equalizes up to homotopy.
\end{theorem}
\begin{proof}
Let $i_0:\mathsf G\to J\to\mathsf F$ be augmentation followed
by the unit. Lemma~\ref{lem:vol2-relative-dual-number-action}
makes $i$ and $i_0$ morphisms from the trivial $\mathsf G$-action
to $(\mathsf F,\alpha)$.

For every receiving pointed map $h$ from $(\mathsf H,\alpha_D)$,
the binary class of $hf$ on $(X\setminus\{q\})\times\{q\}$ is
\begin{equation}\label{eq:vol2-monoid-binary-class}
 [\iota]|_{X\setminus\{q\}}\otimes[h_0]_q
       -[h_0]|_{X\setminus\{q\}}\otimes[\iota]_q.
\end{equation}
The two terms lie in distinct K\"unneth bidegrees. The second
is nonzero by Lemma~\ref{lem:vol2-vertex-singleton} and the unit
identity: $[\iota]_q$ acts as $1_Q$.
Changing the ordered product changes the representative by the
cobar boundary used in Theorem~\ref{thm:vol2-pointed-sphere}.
Lemma~\ref{lem:vol2-binary-homotopy} makes this class invariant
under every pointed geometric homotopy. Thus $hf$ cannot be
homotopic to the map through $J$.

Now suppose a receiving monoid-action map
$g:(\mathsf F,\alpha)\to(Z,\beta)$ equalized $i$ and $i_0$.
Forgetting the action retains a homotopy under $J$. In the
four-factor word $gw$, multiplicativity of $g$ and that
homotopy replace the first and third factors by units.
The second and fourth factors cancel, since $j$ is a monoid map
and the additive $\mathsf H$ has inverse $b\mapsto-b$.
Tensoring the comparison homotopies and applying the monoid unit
coherences therefore gives $gw\simeq\mathsf S\to J\to Z$.
But $w\simeq jf$, and $gj$ receives $\alpha_D$, contradicting
\eqref{eq:vol2-monoid-binary-class}.

A final action would receive $(\mathsf F,\alpha)$ and identify
the two composites from the trivial $\mathsf G$-action, because
its mapping space from that action would be contractible.
The contradiction excludes finality.
\end{proof}

For the polynomial Heisenberg algebra at nonzero level,
$D=\partial_{x_1}$ is a nonzero vertex derivation by
\eqref{eq:vol2-Heisenberg-derivations}. The theorem therefore
excludes a final multiplicative action on this boundary over
$\mathbb C$, independently of the pointed-sphere argument.
For $Q=\mathbb C$, by contrast, the action category is the
slice of transverse monoids over $J$: an action
$F\vot J\to J$ is a monoid map $F\to J$.
Its final object is $1_J$, with evaluation $J\vot J\simeq J$.
The nonzero-derivation hypothesis thus distinguishes an actual
positive case. The ambient endomorphism object remains defined
in both cases, with its different universal property.

\section{Recipient bars and the role of maps}
\label{sec:vol2-recipient-bars}

The density comparison depends on the maps of its index,
not just on the available actors. A weighted bar records
which two actors meet in a common recipient, and how
different choices of their incoming maps are related.

Let $\mathcal T$ be a small infinity-category, let
$Q:\mathcal T^2\to D(k)$ be a coefficient diagram, and
write $\Delta:\mathcal T\to\mathcal T^2$ for the diagonal.
For $d=(a,b)$ the category $d\downarrow\Delta$ consists
of recipients $t$ with maps $a\to t$ and $b\to t$,
including their mapping spaces and coherences. Put
\[
 W(d)=C_*(|d\downarrow\Delta|;k),\qquad
 \widetilde W(d)=\operatorname{fib}(W(d)\to k).
\]
Chains have nonpositive cohomological degrees and the map
is augmentation. An empty recipient category gives the
same formula with $W(d)=0$.

\begin{proposition}\label{prop:vol2-full-recipient-coend}
For the full coefficient diagram, there is an equivalence
\begin{equation}\label{eq:vol2-full-recipient-coend}
 \operatorname{fib}\left(
   \operatorname*{hocolim}_aQ(a,a)
       \longrightarrow\operatorname*{hocolim}_{a,b}Q(a,b)
                  \right)
 \simeq\int^{d\in\mathcal T^2}_{\mathbb L}
                   Q(d)\otimes_k\widetilde W(d).
\end{equation}
No splitting of a vacuum coefficient is required.
\end{proposition}
\begin{proof}
The weight $W(d)$ is the homotopy colimit over $t$ of
the chains on $\operatorname{Map}_{\mathcal T^2}(d,\Delta t)$.
Derived co-Yoneda identifies its coend with
$\operatorname*{hocolim}_tQ(\Delta t)$.
The constant weight $k$ gives
$\operatorname*{hocolim}_dQ(d)$, and its augmentation is
the comparison displayed above. Tensor product and coend
preserve colimits. In a stable category they therefore
preserve finite fibres, proving the formula.
\end{proof}

The same proof works for coefficients in the stable
differential-module category on $U_P$, with its tensoring
over $D(k)$. For the binary density comparison use
\[
 Q(a,b)=(F_{a,I_1}\boxtimes F_{b,I_2})|_{U_P}.
\]
Factorization identifies its diagonal value with $X_P(F_a)$.
Thus \eqref{eq:vol2-full-recipient-coend} computes $K_{C,P}$
on its full index, before any coefficient reduction.

For an ordinary category $\mathcal D$ and strict resolved
diagrams, the normalized two-sided bar has length-$n$ term
\begin{equation}\label{eq:vol2-weighted-bar}
 \bigoplus_{d_0\to\cdots\to d_n}
          \bigl(Q(d_0)\otimes_k\widetilde W(d_n)\bigr)[n].
\end{equation}
Degenerate strings are removed. The first face transports
$Q$, the interior faces compose arrows, and the last pulls
the weight back. The differential is
$\partial+(-1)^nd_{\mathrm{int}}$, with
$\partial=\sum_{i=0}^n(-1)^id_i$ for $n>0$ and no face
term at length zero. Total degree is internal degree
minus $n$, and realization uses direct sums. For an
infinity-category the bar uses chains on mapping spaces,
not just sets of arrows.

If $\mathcal D$ has a terminal object $t$, the bar of
a covariant diagram $V$ contracts to $V(t)$. Let
$\epsilon$ transport length-zero coefficients to $t$
and let $i$ include that summand. On a length-$n$ string,
append the unique arrow to $t$ with sign $(-1)^{n+1}$,
and set the result to zero if that arrow is an identity.
The resulting degree-minus-one map satisfies
\begin{equation}\label{eq:vol2-terminal-contraction}
 dh+hd=1-i\epsilon.
\end{equation}
All faces except the last cancel against the corresponding
faces of the original string. The last gives the identity;
at length zero the other face gives $-i\epsilon$.
The length change reverses the internal differential sign,
so the internal terms cancel too.

\subsection{Equality patterns}

Let $\mathrm{FinInj}$ be a skeleton of finite sets and
injections, with $V(S)=k[S]$. Its monoidal product is
disjoint union and its initial unit is the empty set.
For $r\geq1$, consider
\[
 \operatorname*{hocolim}_S V(S)^{\otimes r}
       \longrightarrow
       \left(\operatorname*{hocolim}_SV(S)\right)^{\otimes r}.
\]
If $B_r$ is the number of partitions of $[r]$, this is
the summation map $k^{B_r}\to k$ in degree zero.
Its fibre is $k^{B_r-1}$.

Indeed the source bar is the linearized nerve of the
category of pairs $(S,t)$ with $t\in S^r$. An injection
preserves the equality partition of $t$. Each such
component has an initial object: the set of blocks,
with its canonical tuple. Its unique injection to any
tuple of that pattern determines every element. Thus
each partition contributes one contractible component.
The case $r=1$ computes the target as $k$, and every
pattern maps to its generator. The discrete partition
gives a section. If all finite-set maps are allowed,
the singleton is terminal, so both sides contract to
$k$ and the comparison is an equivalence.

The first fibre dimensions are $0,1,4,14,51,202,876,4139$.
For two singleton inputs, injections retain separate
same-image and distinct-image recipient components.
A quotient map joins them. Common upper bounds and an
initial unit therefore do not imply diagonal cofinality.

A pointed variant gives the split map explicitly.
Let $\mathsf I$ be finite pointed sets and
point-preserving injections, and put
\[
 V=\operatorname*{hocolim}_{\mathsf I}k[S],\qquad
 W=\operatorname*{hocolim}_{\mathsf I}k[S^2].
\]
The basepoint is a nonzero basis vector. There are two
one-entry patterns, and five two-entry patterns
\[
 (0,0),\quad(0,x),\quad(x,0),\quad(x,x),\quad(x,y),
       \qquad 0,x,y\text{ distinct}.
\]
Their minimal pointed sets are initial in their
components. Hence $V=k^2$ and $W=k^5$, in degree zero.
The comparison $d:W\to V\otimes V$ and its section are
\begin{equation}\label{eq:vol2-pointed-matrices}
 d=\begin{pmatrix}
 1&0&0&0&0\\0&1&0&0&0\\0&0&1&0&0\\0&0&0&1&1
 \end{pmatrix},\qquad
 s=\begin{pmatrix}
 1&0&0&0\\0&1&0&0\\0&0&1&0\\0&0&0&0\\0&0&0&1
 \end{pmatrix}.
\end{equation}
The section comes from Cartesian products of pointed
sets: $(x,y)$ is sent to $((x,0),(0,y))$.
When both entries are nonzero this selects distinct
images. Thus $ds=1$, but the kernel is generated by
$(0,0,0,1,-1)^T$, on which $1-sd$ is the identity.
The cofiber is $k[1]$. This is a computation on the
stated finite-set category, not on $\vAct$.

\subsection{Parallel maps}

Let $\mathcal U$ have two objects $a,b$ and two distinct
nonidentity arrows $f,g:a\to b$. For the constant
diagram $Q=k$ on $\mathcal U^2$, the fibre of its
diagonal comparison has
\[
 H^{-1}=k,\qquad H^0=k,
\]
and no other cohomology. The nerve of $\mathcal U$ is
a circle, and that of $\mathcal U^2$ is a torus.
The diagonal is an isomorphism on $H_0$ and the
diagonal inclusion $k\to k^2$ on $H_1$. Its fibre
therefore receives the target $H_2$ in cohomological
degree $-1$ and the quotient of $H_1$ in degree zero.
Adjoining $h:b\to c$ with $hf=hg$ makes $c$ terminal.
The comparison then contracts to the identity of $k$.

Even the recipient components do not suffice by
themselves: for $(a,a)$ there are three components.
The choices $(f,f)$ and $(g,g)$ connect through
$(a,\vid,\vid)$, whereas $(f,g)$ and $(g,f)$ are
isolated. The coend also retains transport in the
source variables and every bar face. A functional
detecting a density kernel class must be a cocycle
on this whole bar. In degree zero its components
must obey $\ell_bV(f)=\ell_a$ for every arrow;
its higher components obey the dual bar equations.

\section{Evaluation and the remaining homotopies}
\label{sec:vol2-density-homotopies}

The density kernel acts trivially on the boundary.
Its own vanishing is a stronger assertion. We give
the action nullhomotopy first, then write the equations
for a nullhomotopy of the projector itself.

Fix a binary partition $P$. On $U_P$ abbreviate
\[
 Z=X_P(L_C),\quad W=Y_P(L_C),\quad
 A=X_P(C),\quad A'=Y_P(C),\quad
 \delta=\delta_{P,L_C},\quad t=\delta_{P,C}.
\]
The map $t$ is an equivalence. Let $s$ be the density
section and $p=1-s\delta$. The action supplies
$\alpha:Z\vot A\to A$ and
$\alpha^Y:W\vot A'\to A'$.
Choose resolved models of their full diagram. For
homogeneous maps use $\partial u=du-(-1)^{|u|}ud$.
The prescribed paths give degree-minus-one maps
$e,n,h$ and a chain map $g:A'\to A$ with
\[
 \partial e=\delta s-1,\qquad
 \partial n=t\alpha-\alpha^Y(\delta\vot t),\qquad
 \partial h=gt-1.
\]
Put
\begin{equation}\label{eq:vol2-density-evaluation-homotopy}
 \sigma=-e\delta,\qquad
 N=\alpha^Y(\sigma\vot t)+n(p\vot1),\qquad
 H_{\mathrm{ev}}=gN-h\alpha(p\vot1).
\end{equation}
Then $\partial\sigma=\delta p$,
$\partial N=t\alpha(p\vot1)$, and hence
\[
 \partial H_{\mathrm{ev}}=\alpha(p\vot1).
\]
These identities prove the asserted nullhomotopy.
All its paths come from the same enhanced diagram,
so it restricts coherently along the density cocone
and every internal collision. Evaluation on a closed
boundary state, followed by a closed functional,
therefore cannot detect this kernel on cohomology.

\subsection{A map retaining the fibre homotopy}

Set $Z_E=X_P(\vEnd)$, $W_E=Y_P(\vEnd)$ and
$v=\delta_{P,\vEnd}$. Define its factorization defect
$\operatorname{Def}_P(\vEnd)=\operatorname{Fib}(v)$,
using the cochain convention
\begin{equation}\label{eq:vol2-density-fibre-convention}
 \operatorname{Fib}(u)^n=U^n\oplus V^{n-1},\qquad
 d(x,y)=(dx,ux-dy),\qquad u:U\to V.
\end{equation}
Write $\beta_0=X_P(\beta_L)$ and
$\beta_1=Y_P(\beta_L)$. Their coproduct square gives
$\partial\eta=v\beta_0-\beta_1\delta$.
It induces
\[
 T_\eta:K_{C,P}\to\operatorname{Def}_P(\vEnd),
 \qquad T_\eta(x,y)=(\beta_0x,\beta_1y+\eta x).
\]
The map
\begin{equation}\label{eq:vol2-density-detector}
 q_P=T_\eta\pi:Z\to\operatorname{Def}_P(\vEnd),\qquad
 \pi=(p,-e\delta),\qquad
 q_P=(\beta_0p,\eta p-\beta_1e\delta)
\end{equation}
is a cochain map. Indeed $\partial(-e\delta)=\delta p$
is exactly the fibre equation. If $i:K_{C,P}\to Z$
projects to the first component, $i\pi=p$, while
\[
 R(x,y)=(sy,ey),\qquad dR+Rd=1-\pi i.
\]
Thus $\pi i\simeq1$, and $q_P$ is nullhomotopic if
and only if $T_\eta$ is nullhomotopic.

The two components in \eqref{eq:vol2-density-detector}
record an insertion discrepancy. More explicitly,
let $\ell=\beta_0s$ and $b=\eta s+\beta_1e$.
Then $\partial b=v\ell-\beta_1$, and
\[
 (\beta_0-\ell\delta,\eta-b\delta)=q_P.
\]
For an action $a$, restrict this formula along its
density map and inverse factorization, defining
\begin{equation}\label{eq:vol2-insertion-class}
 \Gamma_a=
 [q_PX_P(\lambda_a)\delta_{P,F_a}^{-1}].
\end{equation}
Its lift $\ell Y_P(\lambda_a)$ is precisely the
ordered insertion lift: on the tuple $(a,a)$,
\eqref{eq:vol2-section-formula} uses the two unit
insertions into $a\circledast a$, followed by
composition in $\vEnd$. The path $b$ is their
coproduct and unit path. Composition of homotopy
squares retains the term
$\eta'\beta_0+\beta'_1\eta$, so this identification
includes the second fibre component.

A nonzero $\Gamma_a$ would force $K_{C,P}\ne0$.
Projection to $Z_E$ alone can lose such a class.
Evaluation loses more: its coproduct square sends
$\operatorname{Def}_P(\vEnd)\vot A$ to the fibre
of $t$, which is zero. Conversely, vanishing of
$q_P$ would leave any part of the density kernel
annihilated by $T_\eta$ undetermined.

\subsection{Equations on the full action bar}

Choose a simplicial-category presentation of
$\mathscr T_C$ and a resolved enriched diagram
$X_a=X_P(F_a)$. Let
$Q(a,b)=C_*(\operatorname{Map}_{\mathscr T_C}(a,b);k)$,
in nonpositive cohomological degrees. The constant
right module $k$ uses its augmentation. The two-sided
bar has terms
\begin{equation}\label{eq:vol2-action-bar}
 B_n=\bigoplus_{a_0,\ldots,a_n}
 Q(a_{n-1},a_n)\otimes\cdots\otimes Q(a_0,a_1)
                                      \otimes X_{a_0}.
\end{equation}
Normalize in the simplicial direction, denoting the
result by $NB_n$, and put
\[
 B=\bigoplus_{n\geq0}NB_n[n],\qquad
 D|_{NB_n[n]}=(-1)^nd_{\mathrm{int}}
                        +\sum_{j=0}^n(-1)^jd_j.
\]
The face sum is absent for $n=0$. Simplicial identities
and compatibility with the internal differential give
$D^2=0$. This bar computes $Z$: resolve the constant
weight by represented right modules, whose augmented
bar has the extra identity degeneracy, then tensor
with the coefficient diagram. All tensors are derived.

Let $p_B:B\to B$ represent $p$, and let
$\iota_n:NB_n\to B$ be the inclusion, of degree $-n$.
A degree-minus-one map $H$ with $DH+HD=p_B$ is
equivalent to maps $h_n=H\iota_n$, of degree $-n-1$,
such that
\begin{equation}\label{eq:vol2-full-primitive-equations}
 \partial h_n=p_B\iota_n-
              \sum_{j=0}^n(-1)^jh_{n-1}d_j,
 \qquad
 \partial h_n=Dh_n+(-1)^nh_nd_{\mathrm{int}}.
\end{equation}
The graded Leibniz rule proves the forward implication;
the same equations on every summand define $H$ on
the direct sum and prove the converse.

For an object $a$ and arrow $f:a\to b$, the first
equations are
\begin{equation}\label{eq:vol2-object-arrow-primitives}
 \partial h_a=p_B\iota_a,\qquad
 \partial h_f=p_B\iota_f-h_bX(f)+h_a.
\end{equation}
The cocone term $p_B\iota_f$ remains after the object
equations are solved. A mapping chain $\gamma$ of
degree $-1$, with $d\gamma=g-f$, gives
\[
 \partial\iota_\gamma=\iota_bX(\gamma)-\iota_g+\iota_f,
 \qquad
 \partial h_\gamma=p_B\iota_\gamma-h_bX(\gamma)+h_g-h_f.
\]
Thus paths between action maps impose their own
degree-minus-three primitives.

Unit insertion supplies part of these equations.
Let $D_2a=a\circledast a$ and put
\[
 w_a=\delta_{P,F_{D_2a}}^{-1}
       ((v_1)_{I_1}\boxtimes(v_2)_{I_2})\delta_{P,F_a},
 \qquad r_a=X_P(v_1)-w_a.
\]
The map $w_a$ is a coefficient map; it has not thereby
become an action morphism $a\to D_2a$. Let $V$ and
$W_0$ be the induced maps on the bar from $D_2$ with
$X(v_1)$ and $w$, retaining their coherences. The
section formula gives $W_0\simeq s\delta$; choose
$p_B=1-W_0$. The natural transformation $v_1:1\to D_2$
has a prism $P$ with $DP+PD=V-1$. On an ordinary
string it inserts $a_j\to D_2a_j$ in slot $j$ with
sign $(-1)^j$ and applies $D_2$ thereafter. Interior
faces cancel in pairs and the end faces give $V-1$.
For mapping-space chains the naturality simplices
supply the same cancellations. Consequently
\begin{equation}\label{eq:vol2-unit-residual}
 p_B=\partial(-P)+R_0,\qquad
 R_0=V-W_0,\qquad
 R_0\iota_a=\iota_{D_2a}r_a.
\end{equation}
The projector is nullhomotopic exactly when $R_0$ is.

The residual equations therefore begin with
\begin{equation}\label{eq:vol2-residual-primitives}
 \partial g_a=\iota_{D_2a}r_a,\qquad
 \partial g_f=R_0\iota_f-g_bX(f)+g_a,
\end{equation}
and continue by \eqref{eq:vol2-full-primitive-equations}.
A specified monoid extension supplies an object
primitive: let $q=\mu_M(f\circledast f)$ and let its
unit and action paths give $\partial\tau=X_P(q)r_a$.
Using the extended density cocone if necessary,
\[
 g_a=\iota_M\tau-\iota_qr_a
\]
has the required differential, since
$\partial\iota_q=\iota_MX_P(q)-\iota_{D_2a}$.
The extension category carries its higher coherences.
Extending them to all action maps is an additional
requirement.

\subsection{Objectwise and infinite compatibility}

Consider a category with initial object $j$ and two
parallel arrows $f,g:a\to b$, whose composites from
$j$ agree. Take $V_j=V_b=0$ and $V_a=k$. Its
normalized bar is
\[
 T=[k\{f,g\}\xrightarrow{(-1,-1)}k\{a\}],
\]
in degrees $-1,0$. Every object inclusion is
nullhomotopic, since the one at $a$ sends $1$ to
$-f$ under a primitive. Yet $H^{-1}(T)=k\{f-g\}$,
so $1_T$ is not nullhomotopic. Adding a constant
diagram $k$ gives a vacuum summand. The projector
onto $T$ annihilates that vacuum and every object
inclusion up to homotopy, while preserving the
parallel-arrow class. This finite example isolates
the missing arrow equations.

There is also an inverse-limit condition. Let
$F_mB$ be the subcomplex of bar length at most $m$.
Let $M_m=\vRH(F_mB,B)$ be the derived $k$-linear mapping
complex in the differential-module category on $U_P$.
Since $B=\operatorname*{colim}_mF_mB$,
\[
 \vRH(B,B)\simeq\operatorname*{Rlim}_mM_m.
\]
The fibre model of $1-\mathrm{shift}$ on
$\prod_mM_m$, and exactness of products of vector
spaces, give
\begin{equation}\label{eq:vol2-density-Milnor}
 0\longrightarrow\lim_m{}^1 H^{-1}(M_m)
 \longrightarrow H^0\vRH(B,B)
 \longrightarrow\lim_mH^0(M_m)\longrightarrow0.
\end{equation}
No vanishing of the first term has been proved here.
Nullhomotopies on finite skeleta must be compatible
before they define one on the whole bar.

For the split projector itself there is a stronger criterion.
The finite skeleta in this criterion retain the full coefficient
complexes and every mapping-space chain; only bar length is bounded.

\begin{proposition}\label{prop:vol2-projector-skeleta}
Let $B=\operatorname*{colim}_m B_m$ be an exhaustive filtered
presentation in the derived category of differential modules on
a block-disjoint open. Suppose $B\simeq Y\oplus K$ and
$p\simeq\operatorname{diag}(0,1_K)$.
Then $K=0$ if and only if $H^q(p i_m)=0$ for every $m,q$,
where $i_m:B_m\to B$ is the colimit map.
In particular, $p i_m\simeq0$ for every $m$ implies $p\simeq0$.
\end{proposition}
\begin{proof}
On each affine ordered chart, filtered colimits of differential
modules are exact, so
\[
 H^q(B)=\operatorname*{colim}_mH^q(B_m).
\]
The stated vanishing makes $H^q(p)$ zero on the image of every
stage, hence zero everywhere. But its restriction to $H^q(K)$
is the identity. Thus $K$ is acyclic on every affine chart.
Those restrictions are jointly conservative, so $K=0$.
The converse and the last assertion follow from the splitting.
\end{proof}

Apply the proposition to $p=1-s\delta$ and $B_m=F_mB$.
A nonzero density kernel is therefore detected on some full
finite bar skeleton. No bound on that skeleton follows, and no
finite truncation of states, normal derivatives or mapping chains
is asserted. The proposition also does not make the whole
$\lim^1$ group in \eqref{eq:vol2-density-Milnor} vanish.

Prescribed paths retain a different obstruction. If
$r:M\to N$ restricts the complex of natural transformations
to a specified extension diagram, its fibre has differential
$d(x,y)=(dx,rx-dy)$. A prescribed cancellation $dh=rp$ gives
the cycle $(p,h)$ in this fibre. It need not vanish when the
absolute class $[p]$ vanishes: if $p=dk$, the remaining class
is represented by the closed element $h-rk$ of $N^{-1}$,
modulo the image of $H^{-1}(M)$. This relative cycle is not a
split projector on a coefficient object. The proposition does
not supply its required path.

\begin{problem}\label{qn:vol2-geometric-recipient}
For the Heisenberg boundary $C_{V_\kappa}$ at $\kappa\ne0$,
exhibit a nonzero class of a binary kernel $K_{C,A|B}$ in
the full action bar, with its object, arrow and higher components.
Determine whether such a class can be an insertion class
\eqref{eq:vol2-insertion-class}.
Over $\mathbb C$, determine whether there are two objects of
$\vMon$ with no common receiving multiplicative factorizing
action. The parallel maps of
Theorem~\ref{thm:vol2-monoid-nonfinality} have no receiving
equalizer; this does not decide the two-object question.
\end{problem}

The tree and nilpotent constructions give actual
families of incoming actions. The density section,
evaluation nullhomotopy and unit prism describe their
images and the remaining equations on the full bar.
For this Heisenberg boundary,
Theorem~\ref{thm:vol2-pointed-sphere} proves that these equations
cannot all have a compatible solution: the pointed-action
functor is not representable, and some density kernel is nonzero.
The kernel criterion and primitive equations still apply to
other boundaries. Over $\mathbb C$,
Theorem~\ref{thm:vol2-monoid-nonfinality} also excludes finality
among multiplicative actions, using the multiplication itself.
Common recipients for two specified action objects remain a
distinct comparison problem.

\endgroup
